% Pass options to packages before documentclass to avoid option clash
\PassOptionsToPackage{hyphens}{url}
\PassOptionsToPackage{colorlinks=true,linkcolor=blue,urlcolor=blue,citecolor=black}{hyperref}
\PassOptionsToPackage{authoryear,round}{natbib}

\documentclass[pdflatex,sn-basic]{sn-jnl}% Basic style with author-year citations

% Math packages first
\usepackage{amsmath}%
\usepackage{amssymb}%
\usepackage{amsfonts}%
\usepackage{amsthm}%
\usepackage{amsopn}%
\usepackage{amstext}%
\usepackage{mathrsfs}%
\usepackage{stmaryrd}% For llbracket, rrbracket
\usepackage{mathpartir}%

% Graphics and tables
\usepackage{graphicx}%
\DeclareGraphicsExtensions{.png,.pdf,.jpg}% PNG first, then PDF
\usepackage{subcaption}% For subfigure environment
\usepackage{caption}% For caption customization
\captionsetup{justification=raggedright, singlelinecheck=false}% Left-align all captions by default
\captionsetup[figure]{justification=centering, singlelinecheck=false}% Center-align figure captions
\usepackage{tikz}%
\usetikzlibrary{positioning,fit,backgrounds,calc,arrows.meta,decorations.pathreplacing}%
\usepackage{multirow}%
\usepackage{booktabs}%
\usepackage{pifont}% For \ding{51} (checkmark) and \ding{55} (cross)
\usepackage{tabularx}
\usepackage{array}

% Algorithm packages
\usepackage{algorithm}%
\usepackage{algorithmicx}%
\usepackage{algpseudocode}%

% Other packages
\usepackage{float}% For [H] option to force exact placement
\usepackage{xcolor}%
\usepackage{listings}%
\usepackage{comment}
\usepackage{natbib}% Options passed via \PassOptionsToPackage to avoid clash with sn-jnl
\newcolumntype{L}{>{\raggedright\arraybackslash}X}

% Define Solidity language for listings
\lstdefinelanguage{Solidity}{
  keywords={contract, function, modifier, public, private, external, internal, view, pure, payable, returns, return, if, else, for, while, require, assert, mapping, address, uint256, uint, bool, string, memory, storage, calldata, event, emit, struct, enum, constructor, msg, this, new, delete, break, continue, throw, revert},
  keywordstyle=\color{blue}\bfseries,
  ndkeywords={true, false, null, this, new},
  ndkeywordstyle=\color{darkgray}\bfseries,
  identifierstyle=\color{black},
  sensitive=true,
  comment=[l]{//},
  morecomment=[s]{/*}{*/},
  commentstyle=\color{purple}\ttfamily,
  stringstyle=\color{red}\ttfamily,
  morestring=[b]',
  morestring=[b]"
}

% Listings default settings
\lstset{
  language=Solidity,
  basicstyle=\footnotesize\ttfamily,
  backgroundcolor=\color[RGB]{248,248,248},
  numbers=left,
  numbersep=6pt
}

% Define postbreak symbol for listings
\newcommand{\lstcontinuemark}{\mbox{\textcolor{gray}{\(\hookrightarrow\)}}\space}

\theoremstyle{thmstyleone}%
\newtheorem{theorem}{Theorem}%  meant for continuous numbers
\newtheorem{proposition}[theorem]{Proposition}%
\newtheorem{lemma}[theorem]{Lemma}%
\newtheorem{corollary}[theorem]{Corollary}%
\theoremstyle{thmstyletwo}%
\newtheorem{example}{Example}%
\newtheorem{remark}{Remark}%

\theoremstyle{thmstylethree}%
\newtheorem{definition}{Definition}%

\raggedbottom

% Using author-year style to match submission guidelines
\bibpunct{(}{)}{;}{a}{,}{ } % a = author-year citations

% URL breaking settings (hyperref and url already loaded by sn-jnl.cls)
\Urlmuskip=0mu plus 1mu
\def\UrlBreaks{\do\/\do-\do\_\do\.\do\?\do\&}

\makeatletter
\renewcommand\@biblabel[1]{} % Remove numeric labels like [1]
\makeatother
\setlength{\bibsep}{0.2em}   % Space between bibliography items
\setlength{\bibhang}{2em}

% Custom math commands
\newcommand{\Abort}{\mathsf{Abort}}
\newcommand{\Norm}{\mathsf{Norm}}
\newcommand{\Ret}{\mathsf{Ret}}
\newcommand{\Env}{\mathit{Env}}


\begin{document}

\title[Solidity Intent Checker]{IntentChecker: An Intent-Model-Based Debugging Assistant for Mitigating Numeric Logic Errors in Solidity}

\author[1]{\fnm{Inseong} \sur{Jeon}}\email{iwwyou@korea.ac.kr}

\author[1]{\fnm{Sundeuk} \sur{Kim}}\email{sd\_kim@korea.ac.kr}

\author[1]{\fnm{Hyunwoo} \sur{Kim}}\email{khw0809@korea.ac.kr}

\author*[1]{\fnm{Hoh Peter} \sur{In}}\email{hoh\_in@korea.ac.kr}

\affil*[1]{\orgdiv{Department of Computer Science}, \orgname{Korea University},
\orgaddress{\street{145, Anam-ro}, \city{Seonbuk-gu}, \postcode{02841}, \state{Seoul},
\country{Republic of Korea}}}


\abstract{
Smart contract development lacks the mature debugging infrastructure available in traditional software engineering.
One missing piece is a mechanism for developers to declare their intended numeric behavior and have it checked against actual execution at development time.
For this reason, numeric logic errors---where program behavior diverges from developer intent---are particularly difficult to address.
Existing detection tools instead match completed source code against pre-defined vulnerability patterns in a testing or auditing context.
Errors outside these patterns have no practical path to early prevention.
This paper presents IntentChecker, a development-time intent validation tool for Solidity that operates at multi-contract single-transaction scope and shifts the detection target from pre-defined patterns to developer-defined intent.
At its core is an intent model comprising two annotation types---@During for intended relationships at specific program points within a function, and @Post for intended relationships after function execution---and a semantics that classifies each annotation as Satisfied, Warning, or Violated given the intervals computed by abstract interpretation.
A validation engine applies this intent model to determine whether the developer's intent is violated at each annotation, reporting the result together with a risk score.
We evaluate IntentChecker on real-world numeric logic errors and show that it provides a detection signal complementary to existing pattern-based tools, with analysis latency suitable for development-time use.
}

\keywords{Solidity, Numeric Logic Error, Debugging Assistant, Intent Model, Intent Validation}

\maketitle

\section{Introduction}

The debugging toolchain available to smart contract developers remains far less mature than what traditional software engineering provides.
Our prior work, SolQDebug~\citep{solqdebug}, addressed this immaturity by performing statement-level interval-domain abstract interpretation at development time, eliminating the compile-deploy-trace cycle that traditionally makes Solidity debugging slow.
SolQDebug, however, goes no further than reporting raw intervals---leaving correctness judgment entirely to the developer---and its analysis remains confined to a single contract and a single transaction.

With correctness judgment left to the developer, logic errors---bugs where program behavior diverges from developer intent---are particularly hard to catch, since neither the compiler nor the runtime raises any signal.
A recent empirical study~\citep{defisecurity} confirms this difficulty: practitioners identified logic errors as the most difficult vulnerability type to detect.
In response, several tools~\citep{numscout, gptscan, accountingerrors} explicitly target this class of bugs.
Yet they all take the source code alone as their input, operating in an auditing or testing context.
Moreover, what counts as correct in any given function is set by the developer who wrote it, so by the very definition of a logic error, detecting one requires knowing the developer's intent.
However, current tools take a different approach: they rely on a fixed set of known vulnerabilities chosen by their designers, checking code against a pre-specified set of patterns.
This mismatch is most consequential for numeric computation, which sits at the core of most smart contract code: our analysis~\citep{dappscananalysis} of 101 smart contracts from DAppSCAN-source~\citep{dappscan}, covering 1,023 non-trivial functions, finds numeric types in 82\% of them (Table~\ref{tab:variable_distribution}).
Logic errors are therefore most likely to arise in this space, and addressing such numeric logic errors calls for a fundamentally different kind of analysis: one that takes developer intent as direct input rather than attempting to recover it from code.

Meanwhile, developer preferences favor tools that run at development time. The same empirical study~\citep{defisecurity} reports a clear preference for IDE-integrated tools such as extensions and linters over standalone analyzers, favoring development-time tooling. It also reports a strong prioritization of low false negatives (92.3\%) over low false positives (50\%), matching the profile of sound over-approximating analyses such as abstract interpretation.
These observations together motivate a development-time tool that runs while the function is being written and lets the developer directly declare the numeric intent their contract encodes, rather than relying on pre-authored pattern catalogs to infer it.

To that end, we present IntentChecker, a development-time intent validation tool for Solidity built on SolQDebug and extended to a multi-contract single-transaction scope.
The intent model is built around two annotation types that cover different temporal windows of a function's execution: @During for relationships that should hold at specific program points during execution, and @Post for relationships that should hold at function exit. A tri-valued semantics then maps each annotation to one of Satisfied, Warning, or Violated based on the intervals produced by abstract interpretation.
A validation engine applies this intent model to determine whether the developer's intent is violated at each annotation, reporting the result together with a risk score on a 0--10 scale and the sub-intervals where the violation may occur; the underlying interval comparison primitive is mathematically proved sound against the intent model's set-theoretic specification over all topological configurations of two integer intervals.

We evaluate IntentChecker on 89 real-world numeric logic errors drawn from Web3Bugs~\citep{exploitablebugs} and NumScout~\citep{numscout}.
The evaluation organizes itself around three research questions: the extent to which IntentChecker can surface real-world numeric logic errors through its graded signal (RQ1), the intent model's expressibility ceiling for cases that require bug awareness (RQ2), and how IntentChecker's feedback compares with existing pattern-based detection tools in terms of signal characteristics, latency, and coverage complementarity (RQ3).
Importantly, our positioning with respect to existing tools is complementary rather than competitive: IntentChecker fills a distinct place in the development workflow by enabling intent-driven validation during editing, and is intended to coexist with auditing-time detectors rather than replace them.

\begin{table}[!t]
    \centering
    \caption{Distribution of Variable Types in Smart Contract Functions~\citep{dappscananalysis}}
    \label{tab:variable_distribution}
    \small
    \begin{tabular*}{\columnwidth}{@{\extracolsep{\fill}}lcc}
        \hline
        \textbf{Category} & \textbf{Count} & \textbf{Percentage (\%)} \\
        \hline
        Contracts analyzed & 101 & --- \\
        Functions considered for analysis & 1{,}023 & 100.00 \\
        \hline
        Functions using \textbf{uint} & 798 & 78.01 \\
        Functions using \textbf{int} & 5 & 0.49 \\
        Functions using \textbf{uint} and \textbf{int} & 42 & 4.11 \\
        Functions using other types & 178 & 17.40 \\
        \hline
    \end{tabular*}
\end{table}

\smallskip
\noindent\textbf{Contributions.}
\begin{itemize}
    \item We identify an unfilled gap in current Solidity tooling: the absence of a development-time mechanism for declaring and validating numeric intent, and we show that this gap is not a minor engineering oversight but a structural consequence of how existing numeric logic detection tools are designed (§\ref{sec:background}).
    \item We present an intent model with a formal denotational semantics defining each annotation form as a tri-valued outcome against abstract states, together with a validation engine whose interval comparison primitive is proved sound against the semantics, producing a graded signal (Satisfied / Warning / Violated) with a separately computed risk score and violated region (§\ref{sec:design}).
    \item We extend SolQDebug's analysis scope from single-contract to multi-contract single-transaction, introducing the @IReturn debug annotation to handle interface return values at deployment boundaries (§\ref{sec:prelim}).
    \item We evaluate IntentChecker on 89 real-world numeric logic errors and position it against NumScout and GPTScan along three axes---signal characteristics, per-case latency, and coverage complementarity---demonstrating that the three tools occupy distinct points in the development/auditing timeline rather than forming a single comparison along detection accuracy (§\ref{sec:evaluation}).
\end{itemize}

The remainder of this paper is organized as follows.
Section~\ref{sec:background} motivates the development-time intent validation problem and surveys the structural limitations of existing detection tools.
Section~\ref{sec:prelim} reviews the SolQDebug foundations and describes the scope extensions required for intent validation over multi-contract code.
Section~\ref{sec:design} presents the design of IntentChecker, including its intent model and validation engine.
Section~\ref{sec:evaluation} presents evaluation results across the three research questions.
Section~\ref{sec:discussion} discusses the annotation authoring burden and remaining technical limitations.
Section~\ref{sec:related} reviews related work, and Section~\ref{sec:conclusion} concludes.

\section{Background}\label{sec:background}
\subsection{Numeric Logic Errors in Solidity}

Recent research has identified various numeric-related weaknesses in smart contracts~\citep{numscout,gptscan,flyinointment,exploitablebugs,defiattacks}.
Building on these findings, numeric logic errors can be defined as follows:

\begin{definition}[Numeric Logic Error]
Bugs that allow transactions to complete successfully without compile-time or runtime exceptions,
but violate the developer's intended relationships among numerical values during function execution.
\end{definition}

These errors can manifest in different execution contexts.
Single-transaction errors occur within the execution of a single function call---where computation, function calls, and output generation happen in one atomic sequence.
In contrast, multiple-transaction errors only emerge through sequences of function calls, requiring analysis of inter-function dependencies and contract-level invariants.
Examples of multiple-transaction errors include Risky First Deposit and Price Oracle Manipulation~\citep{gptscan, exploitablebugs}, which manifest only through state changes accumulated across multiple function invocations.

This paper focuses on single-transaction errors---those confined to one atomic execution, though possibly spanning multiple contracts through external calls. This focus is shaped by our broader goal: a development-phase debugging assistant that surfaces logic errors as developers write each function, where the developer's intent is still available as a direct input. Multiple-transaction errors are left to future work.
The representative single-transaction numeric logic error patterns identified in prior work are summarized below:

\begin{itemize}
    \item \textbf{Div In Path}~\citep{numscout}: Division results used in conditions can lead to unexpected control flow due to integer truncation.

    \item \textbf{Operator Order Issue}~\citep{numscout}: Arithmetic order matters: a/b*c $\neq$ a*c/b due to intermediate truncation.

    \item \textbf{Exchange Problem}~\citep{numscout}: Rounding errors in exchange calculations may result in zero output despite valid input, or receiving tokens without payment.

    \item \textbf{Precision Loss Trend}~\citep{numscout}: Integer division's inherent truncation tends to favor one side, potentially misaligning with protocol intent.

    \item \textbf{Minor Amount Retention}~\citep{numscout}: Integer division in distributions leaves remainders unallocated to any recipient, causing permanent lock.

    \item \textbf{Erroneous Accounting}~\citep{gptscan, exploitablebugs}: Incorrect implementation of business model formulas introduces errors that accumulate into substantial losses. Representative sub-patterns include:
    \begin{itemize}
        \item Wrong Interest Rate Order: Balance calculations using outdated or prematurely updated interest rates yield incorrect results.
        \item Wrong Checkpoint Order: Incorrect ordering between state updates and checkpoints leads to reward miscalculation.
    \end{itemize}

    \item \textbf{Inconsistent State Updates}~\citep{exploitablebugs}: A required state variable update is omitted, causing subsequent calculations to use stale values. For example, failing to accrue interest before a balance query produces incorrect results.

    \item \textbf{Greedy Contract}~\citep{flyinointment}: Contract can receive funds but lacks proper logic to withdraw them, causing permanent lock.

\end{itemize}

Despite their diversity, these error patterns share a common root: the actual behavior diverges from what the developer intended.
A straightforward consequence is that any tool that hopes to flag such bugs reliably must reason about the intended behavior, not only the written code.
The next subsection examines existing tools for numeric logic errors in light of this requirement.

\subsection{Structural Constraints of Existing Tools}

Several tools target numeric logic errors.
NumScout~\citep{numscout} catalogs five numerical defect patterns---Div In Path, Operator Order Issue, Exchange Problem, Precision Loss Trend, Minor Amount Retention---and detects them through LLM-pruned symbolic execution.
GPTScan~\citep{gptscan} targets logic vulnerabilities including several accounting-related patterns by combining GPT-based scenario matching with lightweight static analysis~\citep{slither}.
SCType~\citep{accountingerrors} introduces a financial type system over Solidity variables---tuples of $\langle$financial type, scaling factor, token unit, address$\rangle$---and detects accounting errors through type checking.
These tools represent the current state of the art and have meaningfully advanced the field for auditing and testing workflows.
For the development-phase debugging assistant goal articulated in §\ref{sec:background}.1, however, they share two structural constraints rooted in their design.

\smallskip
\noindent\textbf{(1) Analysis input: source code at audit time, with the developer no longer in the loop.}
Existing numeric logic error detection tools take source code as their primary input and are positioned within auditing or testing workflows.
They accept a contract (or a set of contracts), possibly augmented with type annotations or natural-language artifacts, and emit findings as a report.
Because these tools run downstream of authoring, the developer of the code under analysis is no longer in the loop, and they take no input beyond what is listed above.

\smallskip
\noindent\textbf{(2) Detection criteria authored by the tool designer, not the developer.}
Within these input constraints, the developer's intent for the specific function never enters the analysis; each tool instead matches the code against a fixed set of detection criteria authored at design time---pattern catalogs of known buggy shapes~\citep{numscout}, financial type rules~\citep{accountingerrors}, or LLM-driven matching against pre-defined vulnerability scenarios~\citep{gptscan}---derived from prior bugs the tool designer has observed or catalogued.
These criteria describe what counts as buggy code in general, not what counts as a violation of the specific developer's intent for the function being analyzed; a bug is therefore caught only to the extent that it happens to fit a shape the tool designer anticipated.
This distinction is methodological, not quantitative: however rich the criteria, they are chosen by tool designers with no access to that specific intent.

\smallskip
\noindent\textbf{These two constraints are structural, not tool-specific.}
We do not characterize the above constraints as failures of any specific tool; the tools cited above are well-engineered within the design space they chose.
The point is that the design space itself---post-hoc pattern-based detection without explicit intent input---requires intent to be guessed rather than supplied, which is a constraint on what can be expressed as input, not on how well the guessing is implemented.
To step outside this constraint, one must either enlarge the input (ask the developer to supply their intent) or move the tool earlier in time (run it during editing, when the developer is still available to supply that intent).
The approach taken in this paper is to do both.

\subsection{Proposed Overview}\label{sec:proposed-overview}

The two structural constraints above suggest a complementary direction rather than a competing one.
A tool that runs alongside the developer in the editor and lets the developer declare their intent directly can fill the design-space gap without contradicting the strengths of existing detectors.
IntentChecker realizes this direction by addressing each constraint with a dedicated component.

\smallskip
\noindent\textbf{(1) Analysis input: source code at development time, with the developer in the loop.}
IntentChecker is built on SolQDebug~\citep{solqdebug}, our prior work that performs statement-level interval-domain abstract interpretation inside the editor as code is being written.
This places IntentChecker in the same temporal slot as the developer who authors the function under analysis, removing the post-hoc boundary that prevents existing tools from receiving developer input.

\smallskip
\noindent\textbf{(2) Detection criteria authored by the developer, via intent annotations.}
On top of this foundation, IntentChecker introduces an intent model consisting of two annotation types---@During and @Post---that let developers declare numeric relationships at specific program points within a function and at function exit, together with a denotational semantics that maps each annotation to one of Satisfied, Warning, or Violated based on the intervals computed by abstract interpretation.
A validation engine reads these intervals, compares them against the declared intent, and reports the judgment together with a risk score on a 0--10 scale and, when applicable, the sub-intervals in which the intent may be violated.

\smallskip
For these two components to operate on realistic DeFi code, we additionally extend SolQDebug's single-contract single-transaction analysis to a multi-contract single-transaction scope covering inheritance, libraries invoked via using directives and delegatecall, interface-typed variables, and a practically useful arithmetic subset of inline assembly.

\smallskip
This design positions IntentChecker as an awareness tool rather than a detector: it does not pronounce ``this is a bug'' but instead brings the developer's attention to the specific program points where the abstract execution cannot confirm the declared intent.
The next section reviews the SolQDebug foundation on which IntentChecker is built.

\section{SolQDebug and its Scope Extension}\label{sec:prelim}

This section presents the analysis foundation that IntentChecker's intent validation builds upon: a development-time pipeline producing the interval-domain abstract states that the validation engine consumes.
§\ref{sec:prelim:foundations} reviews the parts of SolQDebug~\citep{solqdebug} that IntentChecker relies on, and §\ref{sec:prelim:extension} introduces the scope extensions we add for multi-contract single-transaction analysis.
Readers familiar with SolQDebug may skip §\ref{sec:prelim:foundations} and proceed directly to §\ref{sec:prelim:extension}.

\subsection{SolQDebug}\label{sec:prelim:foundations}

SolQDebug is a development-time debugging assistant for Solidity that provides statement-level feedback through interval-domain abstract interpretation.
It addresses two forms of inefficiency in traditional Solidity debugging: the temporal cost of the compile--deploy--setup--trace cycle in tools such as Remix IDE, and the iterative cost imposed by the non-reversible state model of the EVM.
Both inefficiencies are resolved by performing the entire analysis inside the editor, without contacting any blockchain node.
To keep this in-editor operation responsive as code is being written, SolQDebug extends the Solidity ANTLR grammar with seven specialized entry rules that accept partial code fragments---a lone else branch, a function header without a body, a single statement inside an existing block, and so on---and maintains an incremental control-flow graph so that abstract interpretation can resume from the edit point rather than restart on every edit.
Finally, SolQDebug preserves the analyzer's state before a debug annotation is applied and restores it afterwards, so that ranges can be injected and retracted without invalidating previously computed intervals.

\smallskip
\noindent\textbf{Debug annotations: specifying entry-point values as bounded ranges.}
SolQDebug's analysis takes as input the bounded ranges that the developer supplies for each entry-point value the function reads.
Given these inputs, interval-domain abstract interpretation runs through the function and reports the intervals each variable reaches during execution, giving the developer a concrete view of how the function behaves under the chosen scenario.
Debug annotations are how these inputs are supplied; three families cover the three sources of entry-point values:
\begin{itemize}
    \item @GlobalVar---for EVM execution-context variables such as msg.value, msg.sender, and block.timestamp, whose values are supplied by the enclosing transaction and its block rather than by contract storage or function parameters. The developer declares a representative range reflecting the invocation context under which the function is expected to execute.
    \item @StateVar---for state variables held in contract storage, whose values are determined by prior transactions (e.g., totalSupply, a user's balance, a lending pool's reserves). SolQDebug has no way to know what value these variables hold at the start of the function being debugged; the developer declares a representative range instead.
    \item @LocalVar---for function parameters whose values are supplied by the caller. This is how the developer specifies the test cases under which the function should be analyzed. Variables declared inside the function body are not injectable, because overriding values the analyzer computes from the code itself would break soundness.
\end{itemize}
In all three cases, the developer supplies a bounded range for each entry-point value, and the analyzer proceeds under those ranges.
Values the developer chooses not to constrain remain at $\top$ (``any value'') and propagate through the analysis as such---the result is still sound, just too wide to be informative for debugging.

\smallskip
\noindent\textbf{Range in, range out: no built-in judgment.}
SolQDebug makes no judgment about whether the intervals it computes are ``correct''. It neither compares them against any specification nor flags any as buggy, and their interpretation is left to the developer.
Adding this interpretation layer on top of SolQDebug's ``range in, range out'' foundation is the role of IntentChecker's validation engine (§\ref{sec:design}).

\smallskip
\noindent\textbf{Analysis scope.}
SolQDebug analyzes a single contract for a single transaction.
Inter-contract calls, inheritance, and library invocations are out of this scope; the next subsection describes how IntentChecker extends the analysis to handle the realistic multi-contract code where real numeric logic errors tend to live.

\subsection{Scope Extension for Multi-Contract Single-Transaction Analysis}\label{sec:prelim:extension}

This subsection details the three extensions that realize the multi-contract single-transaction scope introduced in §\ref{sec:proposed-overview}.
The scope admits a single atomic transaction touching several contracts in the same source file or its immediate dependencies, while remaining confined to one execution step; the extensions are (a) a deployment-unit model covering in-scope contract resolution, (b) interface-typed variables and the @IReturn debug annotation, and (c) a practically useful subset of inline assembly.

\smallskip
\noindent\textbf{(a) Deployment unit and in-scope resolution.}
Because IntentChecker debugs one function at a time, during analysis it may step into code that the target function calls---inherited parent functions and library functions invoked via Solidity's using for directive.
For the debug annotations the developer placed on the target contract's state variables to remain valid inside such code, the called code must execute over the same storage as the target.
This holds for inherited parents, whose state variables are compiled into the target's own storage layout, and for libraries invoked via delegatecall, which run under the target's address;
we call this collection---target contract, inherited parents, and storage-sharing libraries---the target's deployment unit, and IntentChecker treats all of it as in-scope.
External contracts---separately deployed contracts invoked through an address or an interface-typed variable---are out of this deployment unit by default, because their storage sits at a different address that the target's debug annotations cannot reach.
This default, however, leaves the analyzer unable to reason about the many DeFi functions that depend on external oracle queries or token balance lookups.
Extension (b) below describes how @IReturn allows a restricted form of data injection across this deployment boundary for view and pure interface calls.

\smallskip
\noindent\textbf{(b) @IReturn: Interface Return Value Debug Annotation.}
Interface-typed variables (e.g., IOracle oracle) present a structural problem that debug annotations on state or local variables cannot solve: the implementation code behind the interface is not present in the source file at all, so SolQDebug's abstract interpreter must return $\top$ for any value derived from an interface call.
This $\top$ then propagates into every downstream computation, defeating the purpose of intent validation for any function that reads an oracle or queries an external token contract.

Our observation is that this situation is semantically identical to what the three existing debug annotations already handle.
In all three cases, the developer supplies a bounded range for a value the function reads whose origin lies outside the analyzer's view.
An interface call's return value is a fourth instance of the same pattern and admits the same mechanism, with the developer supplying a bounded range and the analyzer proceeding under it.
This matches the philosophy of debugging that SolQDebug embodies, where the developer supplies the inputs defining the scenario to inspect and the analyzer reports how the function behaves under them.

Interface calls differ from the other three inputs in one important way---they may modify the external contract's state.
The analyzer cannot see the interface's implementation, so it cannot model any such effect; it can only substitute the developer's range as the return value.
A state-modifying interface call would, in general, mutate the callee's storage, and allowing injection on such calls would let the developer specify the return value but silently omit that side effect, producing a model that does not correspond to any real execution.
We therefore restrict this injection to view and pure interface functions, which Solidity explicitly marks as having no state-modifying side effects, making the developer-supplied range a complete description of the call's observable behavior.

We formalize this intuition as @IReturn, a new debug annotation that lets the developer bind a range to the return value of a view or pure interface function.
Two surface syntaxes are supported:
\begin{itemize}
    \item \textbf{Pattern A (variable-qualified).}
    // @IReturn pool.getPrincipal() = [0, 0] binds the return value of getPrincipal on an interface-typed variable (pool must be declared as some interface type).
    \item \textbf{Pattern B (explicit cast).}
    // @IReturn IERC20(token).balanceOf() = [200, 200] binds the return value of balanceOf when the call is written as an explicit cast of an address variable to an interface type.
\end{itemize}
Both patterns may target structured return values through an access chain, including struct members (vault.delayedStrategyParams().performanceFee) and chained interface calls (engine.mochiProfile().maxCollateralFactor()), so that a single return value can be partially specified and partially left as $\top$.

\smallskip
\noindent\textbf{(c) Arithmetic-focused inline assembly support.}
Solidity's inline assembly (assembly \{ ... \}) is written in Yul and frequently appears in arithmetic-heavy code, particularly in high-precision libraries such as FullMath or ABDKMath.
SolQDebug originally treated any inline assembly block as a single opaque region, leaving all affected variables at $\top$.
IntentChecker implements a dedicated Yul visitor that recognizes the subset of Yul used for arithmetic computation and translates it into the same expression representation used by the rest of the analyzer.
Concretely, the visitor supports variable declarations (let x := expr) and assignments (var := expr) inside a yulBlock, and maps the following Yul built-ins to their Solidity operator counterparts:
\begin{itemize}
    \item \textbf{Arithmetic}: add, sub, mul, div/sdiv, mod/smod, exp.
    \item \textbf{Comparison}: lt/gt/slt/sgt/eq.
    \item \textbf{Bitwise}: and, or, xor.
    \item \textbf{Shifts}: shl, shr, sar.
    \item \textbf{Unary}: not, iszero.
    \item \textbf{Three-argument modular}: mulmod and addmod, which are expanded into $(a \cdot b) \bmod n$ and $(a + b) \bmod n$, respectively.
\end{itemize}
This coverage captures the arithmetic that high-precision libraries typically express in Yul to bypass Solidity's overflow checks, while avoiding the control-flow and memory features (if, for, switch, mload/mstore, sload/sstore) whose abstract semantics would require separate engineering beyond the scope of this paper.
Yul built-ins outside the supported subset (e.g., keccak256, call, log0) are handled conservatively by leaving their results at $\top$, preserving soundness.

\smallskip
\noindent\textbf{Summary.}
Taken together, these three extensions broaden SolQDebug's single-contract single-transaction scope to cover the ways modern Solidity contracts are actually structured: multiple contracts per source file, library-mediated arithmetic, oracle-style interface calls, and hand-written Yul for performance-critical arithmetic.
All extensions preserve SolQDebug's ``range in, range out'' character, so the intent validation layer built on top of this extended analyzer is the subject of the next section.

\section{The Design of IntentChecker}\label{sec:design}
\subsection{System Architecture}

Figure~\ref{fig:architecture} illustrates the overall architecture of IntentChecker.

\begin{figure}[t]
    \centering
    \includegraphics[width=\linewidth]{figure/architecture_drawio.png}
    \caption{System architecture of IntentChecker. Phases~1--2 (Parsing and Analysis Modules) are reused and extended from SolQDebug; Phase~3 (Validation Module) is the core contribution of this work. The Validation Engine receives parsed intent annotations from the Parsing Module and computed intervals from the Analysis Module.}
    \label{fig:architecture}
\end{figure}

\smallskip
\noindent\textbf{Input.}
IntentChecker receives three types of input:
\begin{itemize}
    \item \textbf{Source Fragment}---Solidity source fragments (block or statement) provided incrementally as the developer writes code, as introduced in §\ref{sec:prelim:foundations}.
    \item \textbf{Intent Annotation}---@During and @Post specifications expressing the developer's intended numerical relationships.
    \item \textbf{Debug Annotation}---variable initialization directives that supply bounded value ranges at program points where the actual value is unknown to the analyzer. This category includes SolQDebug's original @GlobalVar (transaction and block context values), @StateVar (state variables in contract storage), and @LocalVar (function parameters), as well as the @IReturn annotation introduced in §\ref{sec:prelim:extension}(b) for interface call return values.
\end{itemize}
In addition, the Analysis Module loads pre-analyzed Dependency CFGs (\texttt{.pkl} files) for libraries and interfaces referenced from the target source file.

\smallskip
\noindent\textbf{Output.}
For each intent annotation, IntentChecker produces:
\begin{itemize}
    \item \textbf{Result}---a validation state of Satisfied, Violated, or Warning.
    \item \textbf{Risk Score}---a 0--10 scale score: $0.0$ for Satisfied, $10.0$ for Violated, and $0.1$--$9.9$ for Warning reflecting the severity of the potential violation.
    \item \textbf{Violated Region}---the specific sub-intervals where the intent may be violated, helping developers identify problematic input ranges.
\end{itemize}

\smallskip
\noindent\textbf{Internal Modules.}
Phases~1 and~2 reuse and extend SolQDebug's analysis pipeline to accommodate the new annotation and scope of §\ref{sec:prelim:extension}, while Phase~3 is IntentChecker's core contribution.
\begin{itemize}
    \item \textbf{Phase~1 --- Parsing Module.} Uses ANTLR for syntax analysis of both Solidity source fragments and intent/debug annotation comments, and additionally delegates the Solidity source fragments to the official Solidity compiler for semantic validation.
    \item \textbf{Phase~2 --- Analysis Module.} Performs incremental abstract interpretation over the interval domain, updating the control-flow graph and variable states as each statement is processed. When a debug annotation is received, the Snapshot Manager (reused unchanged from SolQDebug) preserves the current analysis state, allowing the analyzer to overwrite the targeted variable's interval with the developer-supplied range, re-analyze affected code paths, and restore the pre-annotation state afterwards; the same mechanism handles the new @IReturn annotation, since injecting a range at an interface call return site is mechanically the same as injecting one at a @StateVar declaration.
    \item \textbf{Phase~3 --- Validation Engine.} Draws on two signals from the earlier phases---parsed intent annotations from the Parsing Module, which encode the relational constraints declared in @During/@Post annotations, and computed variable intervals from the Analysis Module, which represent the abstract state of each variable at the annotated program point. For each annotation, it evaluates whether the declared intent holds by comparing the actual intervals against the specified constraints, producing the three output fields listed above.
\end{itemize}

\subsection{Concrete Examples}

This section walks through two examples to illustrate how an annotated function produces a validation judgment.
The first case demonstrates a mid-execution check via a @During annotation; the second demonstrates a post-execution check via @Post together with an @IReturn annotation for external contract state.

\smallskip
\noindent\textbf{Case 1: @During --- Operator Order Issue.}
Figure~\ref{fig:operator-order-example} shows a token purchase function where operator precedence causes the buyer to receive fewer tokens than the exchange rate warrants.

\begin{figure}[t]
\begin{lstlisting}
// contract state: bool presell; uint256 ethRateFix;
// helper: function calculateRate() public view returns (uint256);
function buy() public payable {
    // @GlobalVar msg.value = [50000000000500000, ...]
    // @StateVar presell = true
    require(presell, "presell is closed");
    uint256 amountToBuy = msg.value / ethRateFix * calculateRate();
    // @During amountToBuy * ethRateFix >= msg.value * 250000
    balances[msg.sender] += amountToBuy;
}
\end{lstlisting}
\caption{Operator order issue in exchange rate calculation (simplified).}
\label{fig:operator-order-example}
\end{figure}

The developer's intent is that amountToBuy should satisfy the exchange-rate relation---the tokens received, times the rate denominator, should be at least the ETH paid times the rate numerator.
Because msg.value / ethRateFix truncates before multiplication, this intent is violated for non-round input amounts.
The @During annotation (line~8) captures this intent directly, and IntentChecker reports Violated.

\smallskip
\noindent\textbf{Case 2: @Post --- Erroneous Accounting.}
Figure~\ref{fig:erroneous-accounting-example} shows a rebasing vault's mint function.
The bug is that minted proxy tokens are computed from the total vault balance (including the just-deposited amount), which inflates the shares beyond the deposited amount when the redeem rate deviates from~1.

\begin{figure}[t]
\begin{lstlisting}
// contract state: uint256 _totalSupply; mapping(address => uint256) _balances;
// address baseToken; uint256 ONE; function redeemRate() public view returns (uint256);
function mint(address to, uint256 amount)
    public override returns (uint256) {
    // @LocalVar amount = [500e18, 500e18]
    // @StateVar _totalSupply = [1000e18, 1000e18]
    // @StateVar _balances[to] = [0, 0]
    // @IReturn IERC20(baseToken).balanceOf()
    //   = [1500e18, 1500e18]
    uint256 _redeemRate = redeemRate();
    require(IERC20(baseToken).transferFrom(
        msg.sender, address(this), amount));
    uint256 baseBalance =
        IERC20(baseToken).balanceOf(address(this));
    uint256 proxy = (baseBalance * ONE) / _redeemRate;
    // @Post _balances[to] <= amount
    _mint(to, proxy);
}
\end{lstlisting}
\caption{Erroneous accounting in vault mint with inflated proxy amount (simplified).}
\label{fig:erroneous-accounting-example}
\end{figure}

The @Post \_balances[to] $\leq$ amount annotation (line~16) expresses the intent that no user should receive more shares than the tokens they deposited.
The @IReturn debug annotation (lines~8--9) injects the range $[1500\,\text{e}18, 1500\,\text{e}18]$ for the external balanceOf call (line~14), which the analyzer would otherwise evaluate as $\top$ (§\ref{sec:prelim:extension}(b)).
The transferFrom (line~11) is state-modifying and therefore cannot carry an @IReturn annotation per §\ref{sec:prelim:extension}(b).
With these annotations, IntentChecker computes the proxy amount and finds that it exceeds amount, reporting Violated.
This example demonstrates how @IReturn enables intent validation for functions that depend on external contract state.

\smallskip
These two examples illustrate the two temporal scopes of intent validation---mid-execution checks via @During and post-execution checks via @Post.
The complete intent model supporting these annotations is defined in the following section.

\subsection{Intent Model}\label{sec:intent-model}

The denotational semantics of the intent model is built around a meaning function
\[
\mathcal{V}\llbracket \cdot \rrbracket \;\colon\; \mathbb{A} \;\to\; \bigl(\mathbb{C} \rightharpoonup \mathbb{J}\bigr)
\]
that assigns to each intent annotation $a \in \mathbb{A}$ a partial function from validation contexts to tri-valued judgments. The semantics has three components --- the syntactic domain $\mathbb{A}$, the semantic domain $\mathbb{C}$, and the judgment set $\mathbb{J}$ --- defined as follows.

\smallskip
\noindent\textbf{Syntactic domain.}
$\mathbb{A}$ is the set of intent annotations generated by the grammar of §\ref{sec:intent-semantics}, with a typical element denoted $a$.
Each annotation begins with a marker (@During or @Post) followed by a single body drawn from the corresponding grammatical category.
The grammar partitions bodies into three kinds:
\begin{itemize}
    \item \textbf{@During forms}, attached to an individual statement within the function body, inspect properties at that program point.
    \item \textbf{@Post forms}, evaluated at function exit, inspect function-level properties.
    \item \textbf{Common forms} are shared between the two annotation kinds, applicable in both @During and @Post contexts.
\end{itemize}
There are no compound connectives at the annotation level, so every annotation is atomic.

\smallskip
\noindent\textbf{Semantic domain.}
We use two semantic domains:
\begin{itemize}
    \item $\Sigma$ --- the space of abstract variable environments produced by the interval-domain interpreter. Each $\sigma \in \Sigma$ maps every program variable to an integer interval, or to a compound value (arrays, mappings, or structs) whose leaf components are intervals.
    \item $\mathbb{C}$ --- the set of validation contexts. Each $\Gamma \in \mathbb{C}$ is a record whose fields are listed in Table~\ref{tab:validation-context}; these fields hold the program-state snapshots and auxiliary data the interpreter makes available at an annotated program point, and each annotation form reads only the subset relevant to its form.
\end{itemize}

\begin{table}[t]
\centering
\caption{Fields of the validation context $\Gamma \in \mathbb{C}$. Each atomic annotation form reads only a subset of these fields.}
\label{tab:validation-context}
\small
\begin{tabularx}{\linewidth}{@{} l L @{}}
\toprule
Field & Meaning \\
\midrule
$\sigma_{\text{pt}} \in \Sigma$ & abstract state at the annotated program point \\
$\sigma_{\text{before}} \in \Sigma$ & state immediately before the statement at the annotated line \\
$\sigma_{\text{assign}} \in \Sigma$ & state at a variable's first-assignment site within the enclosing function \\
$\sigma_{\text{entry}} \in \Sigma$ & state at function entry (after parameter binding and debug annotations) \\
$\sigma_{\text{exit}} \in \Sigma$ & joined state at normal function exit \\
$R : \text{Line} \rightharpoonup \text{Val}^*$ & partial map from return-statement lines to captured return values \\
$\mathcal{K}$ & the set of statements at the annotated line (for call-expression and require/assert lookup) \\
$\Pi \subseteq \Sigma$ & the set of abstract states at reachable program points within the enclosing function body, up to the evaluation point \\
\bottomrule
\end{tabularx}
\end{table}

\smallskip
\noindent\textbf{Field supply.}
$\Gamma$ supplies a different subset of fields depending on the annotation kind:
\begin{itemize}
    \item \textbf{@During}: $\sigma_{\text{pt}}$, $\sigma_{\text{before}}$, $\sigma_{\text{assign}}$, $\sigma_{\text{entry}}$, $R$, $\mathcal{K}$, and $\Pi$ (containing every state produced between entry and $\sigma_{\text{pt}}$ along reachable paths);
    \item \textbf{@Post}: $\sigma_{\text{entry}}$, $\sigma_{\text{exit}}$, $R$, and $\Pi$ (containing every state produced along any reachable path from entry to exit).
\end{itemize}

\smallskip
\noindent\textbf{Judgment set.}
$\mathbb{J} = \{S, V, W\}$ is the set of tri-valued judgments assigned to each annotation:
\begin{itemize}
    \item \textbf{Satisfied} ($S$) --- the abstract state guarantees the annotation holds;
    \item \textbf{Violated} ($V$) --- the abstract state guarantees it fails;
    \item \textbf{Warning} ($W$) --- the abstract state admits both satisfying and violating concrete cases.
\end{itemize}
Given an annotation $a \in \mathbb{A}$ whose validation context is $\Gamma \in \mathbb{C}$, IntentChecker reports $\mathcal{V}\llbracket a \rrbracket(\Gamma) \in \mathbb{J}$ as its judgment for that annotation; the partiality of the inner arrow captures the fact that $\mathcal{V}\llbracket a \rrbracket$ is defined only on contexts $\Gamma$ that provide every field $a$ references (see well-formedness in §\ref{sec:intent-semantics}).

The subsection proceeds as follows.
Section~\ref{sec:intent-semantics} fixes the grammar of $\mathbb{A}$ and defines $\mathcal{V}\llbracket \cdot \rrbracket$ via three set-theoretic primitives.
Section~\ref{sec:validation-engine} presents the executable procedures that realise it, with an auxiliary risk score.

\subsubsection{Syntax and Semantics}\label{sec:intent-semantics}

\begin{figure}[t]
\small
\begin{tabular}{@{}r@{\quad}c@{\quad}l@{\qquad}l@{}}
duringIntent & $\rightarrow$ & // @During duringClause & \\[2pt]
postIntent   & $\rightarrow$ & // @Post postClause & \\[6pt]

duringClause & $\rightarrow$ & intentValue (before relOp after)              & $(D_{\text{ba}})$ \\
             & $|$           & intentValue (assign relOp current)            & $(D_{\text{ac}})$ \\
             & $|$           & identifier .arg[ $n$ ] relOp intentValue     & $(D_{\text{arg}})$ \\
             & $|$           & require feasible                              & $(D_{\text{req}})$ \\
             & $|$           & assert feasible                               & $(D_{\text{asrt}})$ \\
             & $|$           & commonClause                                  & $(D_{\text{com}})$ \\[6pt]

postClause & $\rightarrow$ & intentValue (entry relOp exit)                   & $(P_{\text{ee}})$ \\
           & $|$           & commonClause                                     & $(P_{\text{com}})$ \\[6pt]

commonClause & $\rightarrow$ & returnExpression relOp intentValue             & $(C_{\text{ret}})$ \\
             & $|$           & return[$n$] relOp intentValue                  & $(C_{\text{ret}_n})$ \\
             & $|$           & intentValue relOp PercentOf( intentValue , $p$ )  & $(C_{\text{pct}})$ \\
             & $|$           & intentValue relOp ceil( intentValue , $d$ )    & $(C_{\text{ceil}})$ \\
             & $|$           & intentValue relOp floor( intentValue , $d$ )   & $(C_{\text{flr}})$ \\
             & $|$           & intentValue relOp intentValue                   & $(C_{\text{cmp}})$ \\
             & $|$           & intentValue ${=}{>}$ intentValue                & $(C_{\text{impl}})$ \\
             & $|$           & changed ( intentValue , true $|$ false )        & $(C_{\text{chg}})$ \\[6pt]

intentValue & $\rightarrow$ & arithExpr & \\[2pt]
arithExpr   & $\rightarrow$ & arithExpr (+ $|$ -) arithTerm $|$ arithTerm & \\
arithTerm   & $\rightarrow$ & arithTerm (* $|$ / $|$ \%) arithFactor $|$ arithFactor & \\
arithFactor & $\rightarrow$ & number $|$ [ number , number ] $|$ varRef $|$ ( arithExpr ) & \\[2pt]
varRef      & $\rightarrow$ & identifier subAccess* & \\
subAccess   & $\rightarrow$ & . identifier $|$ [ expr ] & \\[6pt]

relOp & $\rightarrow$ & $<$ $|$ $>$ $|$ $\leq$ $|$ $\geq$ $|$ == $|$ != $|$ in $|$ not in & \\
\end{tabular}
\caption{Grammar of intent annotations. Sub-items of $\mathit{duringClause}$, $\mathit{postClause}$, and $\mathit{commonClause}$ are tagged with rule labels $(D_\bullet)$, $(P_\bullet)$, $(C_\bullet)$ for reference in the text.}
\label{fig:intent-grammar}
\end{figure}

\smallskip
\noindent\textbf{Grammar overview.}
Figure~\ref{fig:intent-grammar} gives the BNF of intent annotations, embedded in single-line comments so they coexist with Solidity code without affecting compilation.
Intent-value expressions support arithmetic over variable references with member and index access, numeric literals, inline intervals $[a, b]$, and the DSL helpers PercentOf, ceil, and floor.
Relational operators include $<$, $>$, $\leq$, $\geq$, ==, !=, and the membership operators in and not in; the implication operator ${=}{>}$ is a truthiness-based atomic form whose semantics is defined below.

\smallskip
\noindent\textbf{Well-formedness.}
We write $\Gamma \vdash a : \mathbf{wf}$ to mean that annotation $a$ is well-formed under $\Gamma$: $a$ parses against Fig.~\ref{fig:intent-grammar}, every variable reference in $a$ resolves in the relevant field of $\Gamma$, and all form-specific auxiliary data required by $a$ is present in $\Gamma$.
The equations below define $\mathcal{V}\llbracket a \rrbracket(\Gamma)$ only for well-formed inputs; annotation forms whose auxiliary data is missing correspond to error paths in the implementation and lie outside the semantic claim.

Well-formedness is orthogonal to the informativeness of the input intervals themselves.
Debug annotations (@StateVar, @LocalVar, @IReturn), inherited from SolQDebug, are the mechanism by which developers supply bounded ranges for otherwise-unknown variables; without them, interval inputs may remain $\top$, and the primitives defined below typically return $W$ even though the semantics is still sound.

\smallskip
\noindent\textbf{Expression evaluation.}
Let $\mathcal{I}$ denote the set of integer intervals.
Intent-value expressions are evaluated against a chosen abstract variable environment $\sigma \in \Sigma$ by a function
\[
\text{eval} \;:\; \Sigma \times \mathit{intentValue} \;\to\; \mathcal{I},
\]
which delegates to the interpreter's expression evaluator for ordinary arithmetic subexpressions---inheriting soundness of interval arithmetic from the base analyzer---and to annotation-specific rules for the DSL helpers (PercentOf, ceil, floor) and inline interval literals $[a, b]$.
A numeric literal $n$ evaluates to the singleton interval $[n, n]$, so literals never need to live in $\Gamma$.
By contrast, three rules in Fig.~\ref{fig:intent-grammar} whose left-hand sides are not themselves intent-value expressions are handled outside this function: rule $(D_{\text{arg}})$ is reduced by an auxiliary operator $\text{arg}_n(f, \mathcal{K})$ to the actual argument expression---an intent-value expression---which is then passed to $\text{eval}$, while rules $(C_{\text{ret}})$ and $(C_{\text{ret}_n})$ are resolved directly against the return-summary field $R$ as $R(\ell_{\text{ret}})$ and $R(\ell_{\text{ret}})[n]$, respectively, bypassing $\text{eval}$ because $R$ already stores per-line return intervals.

To support the feasibility forms introduced below, we also use the space of abstract boolean intervals $\mathcal{B} = \{\{0\}, \{1\}, \{0, 1\}\}$ and a boolean-expression evaluator
\[
\text{eval}_{\mathcal{B}} \;:\; \Sigma \times \text{BExpr} \;\to\; \mathcal{B},
\]
where $\text{BExpr}$ is not a non-terminal of Fig.~\ref{fig:intent-grammar} but the base language's (Solidity's) boolean-expression category, from which $\text{cond}_{\text{req}}(\mathcal{K})$ and $\text{cond}_{\text{asrt}}(\mathcal{K})$ extract the boolean condition at the annotated line for rules $(D_{\text{req}})$ and $(D_{\text{asrt}})$, respectively; we inherit the soundness of $\text{eval}_{\mathcal{B}}$ from the underlying analyzer.

\smallskip
\noindent\textbf{Comparison primitive.}
Relational annotations are grounded in the set-theoretic comparison primitive
\[
\text{Cmp}^* \;:\; \mathcal{I} \times \text{Op} \times \mathcal{I} \;\to\; \mathbb{J},
\]
defined as
\[
\text{Cmp}^*(L, \mathit{op}, R) \;=\;
\begin{cases}
S & \text{if } \forall l \in L,\ \forall r \in R :\ l \mathrel{\mathit{op}} r, \\[2pt]
V & \text{if } \forall l \in L,\ \forall r \in R :\ \neg(l \mathrel{\mathit{op}} r), \\[2pt]
W & \text{otherwise.}
\end{cases}
\]
For $\mathit{op} \in \{\text{in}, \text{not\,in}\}$, the set-membership interpretation replaces the pointwise comparison: $l \mathrel{\text{in}} R$ holds iff $l \in R$, and not\,in is its complement.
$\text{Cmp}^*$ gives an algorithm-independent specification of what a relational annotation must evaluate to under a well-formed context.

\smallskip
\noindent\textbf{Truthiness implication primitive.}
The implication operator ${=}{>}$ is an atomic form whose two sides are intent-value expressions rather than full annotations; its meaning is grounded in a truthiness-based set-theoretic primitive
\[
\text{Impl}^* \;:\; \mathcal{I} \times \mathcal{I} \;\to\; \mathbb{J},
\]
defined as
\[
\text{Impl}^*(L, R) \;=\;
\begin{cases}
S & \text{if } L = \{0\}, \\[2pt]
S & \text{if } 0 \notin L \text{ and } 0 \notin R, \\[2pt]
V & \text{if } 0 \notin L \text{ and } R = \{0\}, \\[2pt]
W & \text{otherwise.}
\end{cases}
\]
Intuitively, an interval $X$ is definitely non-zero when $0 \notin X$ and definitely zero when $X = \{0\}$.
$\text{Impl}^*(L, R)$ returns $S$ when $L$ is definitely zero (so the implication holds vacuously) or when both sides are definitely non-zero, $V$ when $L$ is definitely non-zero while $R$ is definitely zero, and $W$ when the zero/non-zero status of either side cannot be established.

\smallskip
\noindent\textbf{Feasibility primitive.}
The feasibility forms are grounded in a two-valued primitive over abstract boolean intervals:
\[
\text{Feasible}^* \;:\; \mathcal{B} \;\to\; \mathbb{J}, \qquad
\text{Feasible}^*(B) \;=\;
\begin{cases}
V & \text{if } B = \{0\}, \\
S & \text{otherwise.}
\end{cases}
\]
$\text{Feasible}^*$ answers whether a Solidity boolean condition---as evaluated by $\text{eval}_{\mathcal{B}}$---can ever hold.
When the abstract interval is $\{0\}$, the condition is provably false on every concrete execution, so the underlying statement is unreachable and rules $(D_{\text{req}})$ and $(D_{\text{asrt}})$ yield $V$.
In every other case the condition admits a true value for at least one input, so the statement is feasible and both rules yield $S$.
Unlike $\text{Cmp}^*$ and $\text{Impl}^*$, $\text{Feasible}^*$ is total but two-valued.
Since the question ``is this condition ever true?'' admits no intermediate answer, the warning case of $\mathbb{J}$ is unused.

\smallskip
\noindent\textbf{Meaning function.}
The meaning function $\mathcal{V}\llbracket \cdot \rrbracket$, whose signature was fixed in §\ref{sec:intent-model}, is defined by case analysis over the annotation grammar, one equation per syntactic form.
We use $e$ (possibly subscripted) for an intent-value expression (in the state-comparison and change-detection forms, typically a variable reference, possibly with struct member or mapping index access), $f$ for a function identifier, and $\ell_{\text{ret}}$ for the return-statement line selected by the annotation: in @During, the return statement on the annotated line; in @Post, the join over all return-statement lines recorded in $R$.

\smallskip
\textbf{@During annotations.}
In the equations below, $\text{arg}_n(f, \mathcal{K})$---serving rule $(D_{\text{arg}})$---denotes the $n$-th argument expression of the call to $f$ located within the statements $\mathcal{K}$ at the annotated line (both sides of the comparison are evaluated in $\sigma_{\text{before}}$ because argument values are fixed immediately before the call executes), and $\text{cond}_{\text{req}}(\mathcal{K})$ and $\text{cond}_{\text{asrt}}(\mathcal{K})$---serving rules $(D_{\text{req}})$ and $(D_{\text{asrt}})$---denote the boolean condition expression of the statement located within $\mathcal{K}$ at the annotated line.
\begin{align*}
\mathcal{V}\llbracket e(\text{before}\ \mathit{op}\ \text{after}) \rrbracket(\Gamma)
    &= \text{Cmp}^*\bigl(\text{eval}(\sigma_{\text{before}}, e),\; \mathit{op},\; \text{eval}(\sigma_{\text{pt}}, e)\bigr), \\[3pt]
\mathcal{V}\llbracket e(\text{assign}\ \mathit{op}\ \text{current}) \rrbracket(\Gamma)
    &= \text{Cmp}^*\bigl(\text{eval}(\sigma_{\text{assign}}, e),\; \mathit{op},\; \text{eval}(\sigma_{\text{pt}}, e)\bigr), \\[3pt]
\mathcal{V}\llbracket f.\text{arg}[n]\ \mathit{op}\ e \rrbracket(\Gamma)
    &= \text{Cmp}^*\bigl(\text{eval}(\sigma_{\text{before}}, \text{arg}_n(f, \mathcal{K})),\; \mathit{op},\; \text{eval}(\sigma_{\text{before}}, e)\bigr), \\[3pt]
\mathcal{V}\llbracket \text{require feasible} \rrbracket(\Gamma)
    &= \text{Feasible}^*\bigl(\text{eval}_{\mathcal{B}}(\sigma_{\text{pt}},\; \text{cond}_{\text{req}}(\mathcal{K}))\bigr), \\[3pt]
\mathcal{V}\llbracket \text{assert feasible} \rrbracket(\Gamma)
    &= \text{Feasible}^*\bigl(\text{eval}_{\mathcal{B}}(\sigma_{\text{pt}},\; \text{cond}_{\text{asrt}}(\mathcal{K}))\bigr),
\end{align*}

\smallskip
\textbf{@Post annotations.}
The only @Post-specific rule $(P_{\text{ee}})$ compares the value of $e$ at function entry ($\sigma_{\text{entry}}$) and exit ($\sigma_{\text{exit}}$); the remaining @Post forms are covered by the common-form equations below.
\begin{align*}
\mathcal{V}\llbracket e(\text{entry}\ \mathit{op}\ \text{exit}) \rrbracket(\Gamma)
    &= \text{Cmp}^*\bigl(\text{eval}(\sigma_{\text{entry}}, e),\; \mathit{op},\; \text{eval}(\sigma_{\text{exit}}, e)\bigr).
\end{align*}

\smallskip
\textbf{Common-form annotations.}
For common-form annotations that may appear in either annotation kind, we let $\text{ref}(\Gamma)$ denote the reference environment: $\text{ref}(\Gamma) = \sigma_{\text{pt}}$ under @During and $\text{ref}(\Gamma) = \sigma_{\text{exit}}$ under @Post, capturing the rule that a common-form annotation is evaluated at the annotated program point inside @During but at the function exit inside @Post.
The following equations apply uniformly to annotations of either kind; rules $(C_{\text{pct}})$, $(C_{\text{ceil}})$, and $(C_{\text{flr}})$ are subsumed by rule $(C_{\text{cmp}})$'s general $e_l\ \mathit{op}\ e_r$ equation below, since their right-hand sides are intent-value expressions that $\text{eval}$ handles as DSL-only constructs:
\begin{align*}
\mathcal{V}\llbracket \text{returnExpression}\ \mathit{op}\ e \rrbracket(\Gamma)
    &= \text{Cmp}^*\bigl(R(\ell_{\text{ret}}),\; \mathit{op},\; \text{eval}(\text{ref}(\Gamma), e)\bigr), \\[3pt]
\mathcal{V}\llbracket \text{return}[n]\ \mathit{op}\ e \rrbracket(\Gamma)
    &= \text{Cmp}^*\bigl(R(\ell_{\text{ret}})[n],\; \mathit{op},\; \text{eval}(\text{ref}(\Gamma), e)\bigr), \\[3pt]
\mathcal{V}\llbracket e_l\ \mathit{op}\ e_r \rrbracket(\Gamma)
    &= \text{Cmp}^*\bigl(\text{eval}(\text{ref}(\Gamma), e_l),\; \mathit{op},\; \text{eval}(\text{ref}(\Gamma), e_r)\bigr), \\[3pt]
\mathcal{V}\llbracket e_l \Rightarrow e_r \rrbracket(\Gamma)
    &= \text{Impl}^*\bigl(\text{eval}(\text{ref}(\Gamma), e_l),\; \text{eval}(\text{ref}(\Gamma), e_r)\bigr).
\end{align*}

The change-detection form $\text{changed}(e, b)$ of rule $(C_{\text{chg}})$ asks whether $e$ has been modified at any reachable program point within the function body relative to its entry value.
Its meaning is therefore not a single $\text{Cmp}^*$ call between two states but a quantification over $\Pi$, lifting $\text{Cmp}^*$ with operator $=$ to the ``exists-a-witness'' reading:
\begin{align*}
\mathcal{V}\llbracket \text{changed}(e, \text{true}) \rrbracket(\Gamma) \;=\;
\begin{cases}
S & \text{if } \exists \sigma \in \Pi : \text{Cmp}^*\bigl(\text{eval}(\sigma_{\text{entry}}, e),\; =,\; \text{eval}(\sigma, e)\bigr) = V, \\[2pt]
V & \text{if } \forall \sigma \in \Pi : \text{Cmp}^*\bigl(\text{eval}(\sigma_{\text{entry}}, e),\; =,\; \text{eval}(\sigma, e)\bigr) = S, \\[2pt]
W & \text{otherwise.}
\end{cases}
\end{align*}
The dual form, asking that $e$ be unchanged throughout the function body, is defined by swapping $S$ and $V$:
\begin{align*}
\mathcal{V}\llbracket \text{changed}(e, \text{false}) \rrbracket(\Gamma) \;=\;
\begin{cases}
S & \text{if } \forall \sigma \in \Pi : \text{Cmp}^*\bigl(\text{eval}(\sigma_{\text{entry}}, e),\; =,\; \text{eval}(\sigma, e)\bigr) = S, \\[2pt]
V & \text{if } \exists \sigma \in \Pi : \text{Cmp}^*\bigl(\text{eval}(\sigma_{\text{entry}}, e),\; =,\; \text{eval}(\sigma, e)\bigr) = V, \\[2pt]
W & \text{otherwise.}
\end{cases}
\end{align*}
Because $\Pi$ is populated differently for @During (states from entry up to the annotated point) and @Post (states from entry up to any exit), these two equations cover both annotation kinds with a single pair of cases; no separate equations per kind are needed.

\subsubsection{Auxiliary Outputs: Violated Region and Risk Score}\label{sec:validation-engine}

\begin{algorithm}[!t]
\caption{Auxiliary Outputs Computation for $\text{Cmp}^*$-based Annotations}
\label{alg:risk-score}
\begin{algorithmic}[1]
\State \textbf{Input:} Intervals $L$, $R$, operator $op$, state $j \in \mathbb{J}$
\State \textbf{Output:} $(\mathit{risk\_score},\; \mathit{regions})$ with $\mathit{risk\_score} \in [0.0, 10.0]$ and $\mathit{regions} = (L_{\mathit{lower}}, L_{\mathit{upper}}, L_{\mathit{overlap}})$
\State $(|T_s|, |T_v|, |T_w|) \gets \mathit{RegionSizes}(L, R, op)$
\State $\mathit{regions} \gets \mathit{ExtractViolatedRegion}(L, R, op)$
\If{$j = \mathit{satisfied}$}
    \State \Return $(0.0,\; \mathit{regions})$
\ElsIf{$j = \mathit{violated}$}
    \State \Return $(10.0,\; \mathit{regions})$
\EndIf
\State $p \gets p_{\mathit{true}}(L, R, op)$
\State $\mathit{risk\_type} \gets \mathit{ClassifyRisk}(\mathit{regions},\, op)$
\State $(\mathit{base},\, \mathit{span}) \gets \mathit{BandOf}(\mathit{risk\_type})$
\State $\mathit{risk\_score} \gets \mathit{base} + (1 - p) \times \mathit{span}$
\State \Return $(\mathit{risk\_score},\; \mathit{regions})$
\end{algorithmic}
\end{algorithm}

This subsection presents the executable procedures that accompany the semantic realisation of §\ref{sec:intent-semantics}, producing two auxiliary outputs derived from the same interval analysis that drives $\text{Cmp}^*$:
\begin{itemize}
    \item \textbf{violated region} --- sub-intervals of the input range $L$ that witness the failure of the comparison;
    \item \textbf{risk score} --- a value on a 0--10 scale combining the state, the true probability, and the shape of the violated region.
\end{itemize}

Both outputs are downstream of the semantic claim and do not refine the judgment $j \in \mathbb{J}$; they support warning ranking and developer feedback.
An internal probability $p_{\mathit{true}}(L, R, op) \in [0, 1]$ --- the chance that the comparison holds under a uniform distribution over $L \times R$ --- mediates the risk-score computation but is not surfaced separately.
$\text{Impl}^*$ and $\text{Feasible}^*$ are directly computable from their definitions and need no dedicated algorithm, and the change-detection form iterates $\Pi$ directly.
The remaining auxiliary outputs for $\text{Cmp}^*$-based annotations are computed by Algorithm~\ref{alg:risk-score} in four steps, each elaborated in turn below and cited with its line range in the algorithm.

\smallskip
\noindent\textbf{Step 1: Region partitioning (Algorithm~\ref{alg:risk-score}, line 3).}
Given integer intervals $L = [l_{\min}, l_{\max}]$ and $R = [r_{\min}, r_{\max}]$ with an operator $op$, the interpreter partitions $L$ into three disjoint regions:
\begin{itemize}
    \item $T_s$ (satisfied region): values of $L$ for which $l \mathrel{op} r$ holds for all $r \in R$;
    \item $T_v$ (violated region): values of $L$ for which $l \mathrel{op} r$ fails for all $r \in R$;
    \item $T_w$ (warning region): values of $L$ whose truth value depends on the actual value of $R$.
\end{itemize}

Table~\ref{tab:op-formulas} gives the satisfied/violated conditions and region sizes per operator, where $n_L = l_{\max} - l_{\min} + 1$ and $n_R = r_{\max} - r_{\min} + 1$ denote the number of integer values in each interval.
The state $j$ returned by $\text{Cmp}^*$ follows directly: $j = S$ when the satisfied condition holds, $j = V$ when the violated condition holds, and $j = W$ when neither holds.
Figure~\ref{fig:region-partition} illustrates this partition for the canonical warning configuration of $op = \text{``$>$''}$.

\begin{table}[t]
\centering
\caption{Per-operator satisfied/violated conditions and region sizes. All sizes are clamped to $\max(0, \cdot)$.}
\label{tab:op-formulas}
\small
\begin{tabular}{cllcc}
\toprule
$op$ & Satisfied condition & Violated condition & $|T_s|$ & $|T_v|$ \\
\midrule
$>$    & $l_{\min} > r_{\max}$    & $l_{\max} \leq r_{\min}$  & $l_{\max} - r_{\max}$     & $r_{\min} - l_{\min} + 1$ \\
$\geq$ & $l_{\min} \geq r_{\max}$ & $l_{\max} < r_{\min}$     & $l_{\max} - r_{\max} + 1$ & $r_{\min} - l_{\min}$     \\
$<$    & $l_{\max} < r_{\min}$    & $l_{\min} \geq r_{\max}$  & $r_{\min} - l_{\min}$     & $l_{\max} - r_{\max} + 1$ \\
$\leq$ & $l_{\max} \leq r_{\min}$ & $l_{\min} > r_{\max}$     & $r_{\min} - l_{\min} + 1$ & $l_{\max} - r_{\max}$     \\
\midrule
$=$    & $L = R$, both singletons  & $L \cap R = \emptyset$  & $|L \cap R|$ & $n_L - |L \cap R|$ \\
$\neq$ & $L \cap R = \emptyset$  & $L = R$, both singletons  & $n_L - |L \cap R|$ & $|L \cap R|$ \\
\midrule
in     & $L \subseteq R$ & $L \cap R = \emptyset$ & $|L \cap R|$       & $n_L - |L \cap R|$ \\
not\,in & $L \cap R = \emptyset$ & $L \subseteq R$ & $n_L - |L \cap R|$ & $|L \cap R|$       \\
\bottomrule
\end{tabular}
\end{table}

\begin{figure}[t]
\centering
\vspace{4pt}
\begin{tikzpicture}[
    x=0.85cm, y=0.5cm,
    Rbox/.style={fill=red!30, draw=red!60!black, line width=0.5pt},
    Tsbox/.style={fill=green!40, draw=green!60!black, line width=0.5pt},
    Tvbox/.style={fill=red!50, draw=red!65!black, line width=0.5pt},
    Twbox/.style={fill=yellow!60, draw=orange!65!black, line width=0.5pt},
    ax/.style={->, line width=0.6pt, black!70},
    tick/.style={line width=0.5pt, black!70},
    txt/.style={font=\scriptsize},
    tny/.style={font=\tiny},
    Tlbl/.style={font=\small},
    braced/.style={decorate, decoration={brace, amplitude=3pt, mirror, raise=1pt}},
]

% R (same thickness as L, above L)
\fill[Rbox] (6, 6.0) rectangle (12, 7.4);
\node[Tlbl] at (9, 6.7) {$R$};
\draw[tick] (6, 5.8) -- (6, 7.6);
\draw[tick] (12, 5.8) -- (12, 7.6);
\node[tny, anchor=south] at (6, 7.6) {$r_{\min}$};
\node[tny, anchor=south] at (12, 7.6) {$r_{\max}$};

% L (thicker, partitioned)
\fill[Tvbox] (2, 3.8) rectangle (6, 5.2);
\fill[Twbox] (6, 3.8) rectangle (12, 5.2);
\fill[Tsbox] (12, 3.8) rectangle (16, 5.2);
\node[Tlbl] at (4, 4.5) {$T_v$};
\node[Tlbl] at (9, 4.5) {$T_w$};
\node[Tlbl] at (14, 4.5) {$T_s$};
\draw[tick] (2, 3.6) -- (2, 5.4);
\draw[tick] (16, 3.6) -- (16, 5.4);
\node[tny, anchor=north] at (2, 3.6) {$l_{\min}$};
\node[tny, anchor=north] at (16, 3.6) {$l_{\max}$};
\node[txt, anchor=east] at (1.8, 4.5) {$L$};

% Braces below L
\draw[braced] (2, 2.8) -- (6, 2.8);
\node[tny, anchor=north] at (4, 2.7) {$|T_v|{=}r_{\min}{-}l_{\min}{+}1$};
\draw[braced] (6, 2.8) -- (12, 2.8);
\node[tny, anchor=north] at (9, 2.7) {$|T_w|{=}r_{\max}{-}r_{\min}$};
\draw[braced] (12, 2.8) -- (16, 2.8);
\node[tny, anchor=north] at (14, 2.7) {$|T_s|{=}l_{\max}{-}r_{\max}$};

% Axis
\draw[ax] (0, 1.8) -- (17, 1.8);
\node[txt, anchor=west] at (17, 1.8) {$\mathbb{Z}$};

\end{tikzpicture}

\vspace{6pt}
\caption{Partition of $L$ on the integer line $\mathbb{Z}$ for $op = \text{``$>$''}$ in the canonical warning case ($l_{\min} < r_{\min} < r_{\max} < l_{\max}$). Region sizes match Table~\ref{tab:op-formulas}'s $op = \text{``$>$''}$ row.}
\label{fig:region-partition}
\vspace{4pt}
\end{figure}

\smallskip
\noindent\textbf{Step 2: Violated-region extraction (Algorithm~\ref{alg:risk-score}, line 4).}
The interpreter also extracts the violated region of $L$: sub-intervals of $L$ that witness failure.
Table~\ref{tab:false-regions} decomposes $L$ into three directional components per operator:
\begin{itemize}
    \item $L_{\mathit{lower}}$ --- values of $L$ that fall below the acceptable range of $R$;
    \item $L_{\mathit{upper}}$ --- values of $L$ that exceed the acceptable range of $R$;
    \item $L_{\mathit{overlap}}$ --- values of $L$ whose outcome depends on the concrete $R$-value (the warning region for ordering operators; the coincidence zone for equality/membership).
\end{itemize}

Each component is an integer sub-interval or empty; empty components indicate that no witness exists in that direction.
For the $=$ operator, when $L$ and $R$ are disjoint, the overlap is empty and $L$ falls entirely into exactly one of $L_{\mathit{lower}}$ or $L_{\mathit{upper}}$ depending on which side it lies on.
The risk-classification step below consumes this decomposition, and the same sub-intervals are surfaced to the developer as the concrete ranges where the annotation may fail.

\begin{table}[t]
\centering
\caption{Violated-region extraction formulas. Endpoints are clamped so that $\min > \max$ yields empty; ``---'' denotes no witness in that direction.}
\label{tab:false-regions}
\footnotesize
\begin{tabular*}{\linewidth}{@{\extracolsep{\fill}} c l l l @{}}
\toprule
$op$ & $L_{\mathit{lower}}$ & $L_{\mathit{upper}}$ & $L_{\mathit{overlap}}$ \\
\midrule
$>$    & $[l_{\min}, \min(l_{\max}, r_{\min})]$         & ---                                             & $[\max(l_{\min}, r_{\min}{+}1), \min(l_{\max}, r_{\max})]$ \\
$\geq$ & $[l_{\min}, \min(l_{\max}, r_{\min}{-}1)]$     & ---                                             & $[\max(l_{\min}, r_{\min}), \min(l_{\max}, r_{\max})]$ \\
$<$    & ---                                            & $[\max(l_{\min}, r_{\max}), l_{\max}]$          & $[\max(l_{\min}, r_{\min}), \min(l_{\max}, r_{\max}{-}1)]$ \\
$\leq$ & ---                                            & $[\max(l_{\min}, r_{\max}{+}1), l_{\max}]$      & $[\max(l_{\min}, r_{\min}), \min(l_{\max}, r_{\max})]$ \\
\midrule
$=$    & $[l_{\min}, (L \cap R)_{\min}{-}1]$            & $[(L \cap R)_{\max}{+}1, l_{\max}]$             & $L \cap R$ \\
$\neq$ & ---                                            & ---                                             & $L \cap R$ \\
\midrule
\text{in}      & $[l_{\min}, \min(l_{\max}, r_{\min}{-}1)]$ & $[\max(l_{\min}, r_{\max}{+}1), l_{\max}]$  & --- \\
\text{not\,in} & ---                                    & ---                                             & $L \cap R$ \\
\bottomrule
\end{tabular*}
\end{table}

\smallskip
\noindent\textbf{Step 3: State-based scoring (Algorithm~\ref{alg:risk-score}, lines 5--9).}
Lines 5--9 map the state $j$ returned by $\text{Cmp}^*$ directly to extreme scores: $j = \mathit{satisfied}$ returns $0.0$ (fully satisfied condition) and $j = \mathit{violated}$ returns $10.0$ (definite violation).
The warning case $j = W$ requires further composition and is handled in Step~4.

\smallskip
\noindent\textbf{Step 4: Warning-case scoring (Algorithm~\ref{alg:risk-score}, lines 10--13).}
When $j = W$, the risk score is composed from an internal probability $p_{\mathit{true}}$ and a risk-type classification over the violated-region components, then assembled into a value in $(0.0, 10.0)$ by a banded linear formula.

\smallskip
\noindent\textit{True probability (Algorithm~\ref{alg:risk-score}, line 10).}
The true probability $p_{\mathit{true}}(L, R, op) \in [0, 1]$ that the comparison holds under a uniform distribution over $L \times R$ is defined in closed form using the region sizes from Step~1:
\[
p_{\mathit{true}}(L, R, op) \;=\;
\begin{cases}
|T_s| / n_L & \text{if } op \in \{\text{in}, \text{not\,in}\}, \\[3pt]
|T_s| / (n_L \cdot n_R) & \text{if } op = \text{``=''} \text{ and } n_R > 1, \\[3pt]
1 - |T_v| / (n_L \cdot n_R) & \text{if } op = \text{``$\neq$''} \text{ and } n_R > 1, \\[3pt]
\bigl(|T_s| + E_w\bigr) / n_L & \text{otherwise.}
\end{cases}
\]

\begin{figure}[t]
\centering
\vspace{4pt}
\begin{tikzpicture}[
    x=0.85cm, y=0.5cm,
    bar/.style={fill=orange!55, draw=orange!75!black, line width=0.4pt},
    ax/.style={->, line width=0.6pt, black!70},
    txt/.style={font=\scriptsize},
    tny/.style={font=\tiny},
]

\foreach \x/\h in {4/0.4, 5/0.8, 6/1.2, 7/1.6, 8/2.0} {
  \fill[bar] (\x-0.35, 0) rectangle (\x+0.35, \h);
}

\draw[ax] (3, 0) -- (9.5, 0);
\draw[ax] (3.3, -0.2) -- (3.3, 2.3);

\node[tny, anchor=north] at (4, -0.1) {$r_{\min}{+}1$};
\node[tny, anchor=north] at (8, -0.1) {$r_{\max}$};
\node[tny, anchor=east] at (3.2, 0.35) {$a$};
\node[tny, anchor=east] at (3.2, 1.95) {$b$};

\node[txt, anchor=west] at (11, 1.0)
  {$\displaystyle E_w = \frac{(a+b)\,|T_w|}{2\,n_R}$};

\end{tikzpicture}

\vspace{6pt}
\caption{Count of $r \in R$ satisfying $l > r$ across $T_w$ for the canonical warning configuration. The count grows arithmetically from $a = 1$ at $l = r_{\min}{+}1$ to $b = n_R - 1$ at $l = r_{\max}$; summing over $T_w$ and dividing by $n_R$ yields $E_w$.}
\label{fig:arithmetic-series}
\vspace{4pt}
\end{figure}

\begin{description}
\item[$op \in \{\text{in}, \text{not\,in}\}$.] $|T_s|$ already counts the $L$-values inside $R$, yielding $|T_s|/n_L$ directly.
\item[$op = \text{``=''}$ or $\text{``$\neq$''}$ with $n_R > 1$.] Each overlapping $L$-value has a $1/n_R$ probability of matching any particular $R$-value, so satisfying pairs are counted over the product space and divided by $n_L \cdot n_R$ (all pairs).
\item[Otherwise ($>$, $\geq$, $<$, $\leq$, and $=$/$\neq$ with $n_R = 1$).] The warning region $T_w$ is the range of $L$-values whose outcome depends on $R$.
Since Solidity operates on integers, for each $l \in T_w$ the number of $R$-values satisfying $l \mathrel{op} r$ forms an arithmetic sequence across $T_w$ (visualised in Figure~\ref{fig:arithmetic-series}), so
\[
E_w \;=\; \frac{(a + b) \, |T_w|}{2 \, n_R},
\]
where $a$ and $b$ are the numbers of satisfying $R$-values at the smallest and largest elements of $T_w$, respectively.
When $n_R = 1$ for $=$ or $\neq$, $E_w = 0$ and the formula degenerates to pointwise evaluation.
\end{description}

\smallskip
\noindent\textit{Risk type classification (Algorithm~\ref{alg:risk-score}, line 11).}
For warning results, the engine classifies the violation pattern into three risk types based on which components of Table~\ref{tab:false-regions} are non-empty:
\begin{itemize}
    \item \textbf{Type 1} --- no directional witness in $L_{\mathit{lower}}$ or $L_{\mathit{upper}}$ (only an overlap zone, representing a mild uncertainty);
    \item \textbf{Type 2} --- exactly one of $L_{\mathit{lower}}$ or $L_{\mathit{upper}}$ is non-empty (a single-direction violation);
    \item \textbf{Type 3} --- both $L_{\mathit{lower}}$ and $L_{\mathit{upper}}$ are non-empty, or the operator is $=$/$\neq$ (whose interval semantics is inherently imprecise and is assigned the highest band by convention).
\end{itemize}

\smallskip
\noindent\textit{Banded scoring (Algorithm~\ref{alg:risk-score}, lines 12--13).}
Each type maps to a fixed $(\mathit{base}, \mathit{span})$ band --- $(0.1, 3.2)$ for type~1, $(3.4, 3.2)$ for type~2, $(6.7, 3.2)$ for type~3 --- and the final score within the band is $\mathit{base} + (1 - p_{\mathit{true}}) \times \mathit{span}$, so that a lower $p_{\mathit{true}}$ (a more probable violation) yields a higher score.

This completes the design of IntentChecker's intent model and validation engine. The full implementation, together with the evaluation benchmark and scripts used in the following section, is publicly available~\citep{intentchecker}.

\section{Evaluation}\label{sec:evaluation}

We evaluate IntentChecker through three research questions that probe its empirical mitigation capability, the structural reasons behind unmitigated cases, and its position relative to existing detection tools:

\begin{itemize}
    \item \textbf{RQ1 (Mitigation Capability)}: Can intent-based validation provide developers with actionable diagnostic signals for numeric logic errors during development?
    \item \textbf{RQ2 (Unmitigated Case Analysis)}: For instances that the tool cannot handle, what is the structural reason, and how does this distribute across limitation types?
    \item \textbf{RQ3 (Comparison with Existing Detection Tools)}: How does IntentChecker's approach compare to existing detection tools in terms of detection model, coverage, and analysis latency?
\end{itemize}

\subsection{Experimental Setup}\label{sec:setup}

\noindent\textbf{Experimental Setting.}
The evaluation environment consists of an 11th Gen Intel\textregistered{} Core\texttrademark{}
i7-11390H CPU at 3.40\,GHz with 16.0\,GB RAM, running Windows 10 (64-bit).

\smallskip
\noindent\textbf{Dataset Collection.}
We construct a benchmark of 89 real-world numeric logic error instances drawn from three complementary sources, each contributing different error patterns under IntentChecker's scope (multi-contract / single-transaction).

\smallskip\noindent\textbf{NumScout (7 instances).}
NumScout~\citep{numscout} provides a curated collection of numeric logic errors mined from deployed BNB Smart Chain contracts.
We use the 95-sample subset distributed in NumScout's official repository~\citep{numscout_dataset}, which the authors identify as their main case-level evaluation set.
We select seven cases covering all five numerical defect patterns enumerated in NumScout (Div In Path, Operator Order Issue, Exchange Problem, Precision Loss Trend, and Minor Amount Retention), with two instances for Precision Loss Trend and two for Exchange Problem---one of which is labeled Profit Opportunity in the repository.
All selected contracts target Solidity~$\geq$0.6, the minimum version supported by IntentChecker's front-end; one contract carries a 0.4.x pragma but uses only constructs semantically identical under $\geq$0.6.

\smallskip\noindent\textbf{Flyinointment (1 instance).}
We additionally include one Greedy Contract instance from the dataset of \citet{flyinointment}, which originally introduced the Greedy Contract pattern and is referenced in §\ref{sec:background}.1.

\smallskip\noindent\textbf{Web3Bugs (81 instances).}
Web3Bugs~\citep{exploitablebugs} systematically collected and classified high-severity vulnerabilities from Code4rena audit contests.
We select bugs whose execution context fits IntentChecker's scope, namely multi-contract single-transaction numeric logic errors.
We draw from (i) categories S6-1 through S6-4 (Erroneous Accounting), whose root causes are semantic errors in arithmetic or accounting logic, and (ii) category S3-1 (Inconsistent State Updates), retaining only the subset whose root cause involves a numeric computation with a missing or stale state variable update.
The resulting Web3Bugs subset contains 81 instances spanning Erroneous Accounting (72) and Inconsistent State Updates (9).

\begin{table}[t]
\centering
\caption{Benchmark dataset composition.}
\label{tab:dataset}
\begin{tabular*}{\linewidth}{@{\extracolsep{\fill}} l l r @{}}
\toprule
\textbf{Source} & \textbf{Error Pattern} & \textbf{Count} \\
\midrule
\multirow{5}{*}{NumScout}
  & Div In Path              & 1  \\
  & Exchange Problem         & 2  \\
  & Operator Order Issue     & 1  \\
  & Precision Loss Trend     & 2  \\
  & Minor Amount Retention   & 1  \\
\midrule
Flyinointment & Greedy Contract & 1 \\
\midrule
\multirow{2}{*}{Web3Bugs}
  & Erroneous Accounting        & 72 \\
  & Inconsistent State Updates  & 9  \\
\midrule
\textbf{Total (collected)} & & \textbf{89} \\
Excluded (see below)        & & $-14$ \\
\textbf{Effective}          & & \textbf{75} \\
\bottomrule
\end{tabular*}
\end{table}

\smallskip\noindent\textbf{Exclusions.}
Of the 89 collected instances, 14 are excluded after manual inspection because they fall outside the definition of a development-time numeric logic error.
The exclusion categories are:
\begin{itemize}
  \item \textbf{E1 (Source already patched, 2).} The provided source code is a post-fix version; the reported bug cannot be reproduced.
  \item \textbf{E2 (Integer overflow/underflow, 2).} The program reverts via Solidity~$\geq$0.8 checked arithmetic rather than returning a wrong value, so no numeric logic error is observable at the output level.
  \item \textbf{E3 (Duplicate, 2).} Duplicate audit findings on the same buggy site.
  \item \textbf{E4 (Not a bug, 1).} The project sponsor disputed the finding as intended behavior.
  \item \textbf{E5 (Missing dependency, 2).} Required external libraries not provided in the audit package.
  \item \textbf{E6 (Multi-transaction, 3).} Outside the scope of this study: the bug requires state changes across separate transactions, which IntentChecker's single-transaction analysis cannot cover.
  \item \textbf{E7 (Language semantics, 1).} The reported issue is an inherent property of Solidity's integer-only arithmetic (truncating division), not a developer error; no alternative implementation can produce a different result.
  \item \textbf{Bug not in target contract (1).} The bug actually resides in a different contract than the one documented in the audit report.
\end{itemize}
The remaining \textbf{75 instances} form the effective evaluation set used in RQ1--RQ3.

\subsection{RQ1: Mitigation Capability}\label{sec:rq1}

\begin{table}[t]
\centering
\caption{Evaluation outcome for the 75 effective benchmark instances. Top: the 20 mitigated cases by annotation form. Bottom: the 55 not\_detectable cases by limitation type (sub-categories defined in §\ref{sec:rq2}).}
\label{tab:rq1-overview}
\small
\begin{tabular*}{\linewidth}{@{\extracolsep{\fill}} l l r @{}}
\toprule
\textbf{Outcome} & \textbf{Detail} & \textbf{Cases} \\
\midrule
\multicolumn{3}{@{}l}{\textbf{Mitigated (20)}} \\
\addlinespace
\multicolumn{3}{@{}l}{\quad @During (12)} \\
& Relational comparison (var relOp expr)        & 5 \\
& Before/After (var(Before relOp After))        & 2 \\
& Function-arg check (func.arg[N] relOp expr)  & 3 \\
& Implication form                              & 1 \\
& Feasibility check (require feasible)          & 1 \\
\addlinespace
\multicolumn{3}{@{}l}{\quad @Post (8)} \\
& Relational comparison (var relOp expr)        & 4 \\
& Entry/Exit (var(Entry relOp Exit))            & 2 \\
& Return expression (returnExpression relOp expr) & 1 \\
& changed(var, true/false)                      & 1 \\
\midrule
\multicolumn{3}{@{}l}{\textbf{Not detectable (55)}} \\
\addlinespace
& L1: Loop-related precision loss               & 10 \\
& L2: Interface / external call returns $\top$  & 4 \\
& L3: Unsupported construct                     & 7 \\
& L4: Inexpressible intent                      & 20 \\
& L5: Bug-awareness required                    & 14 \\
\midrule
\multicolumn{2}{@{}l}{\textbf{Total}} & \textbf{75} \\
\bottomrule
\end{tabular*}
\end{table}

\begin{table}[t]
\centering
\caption{Per-case analysis time on the 20 mitigated cases (mean over 3 runs, single thread, dependency pre-analysis excluded). Cases sorted by mean time.}
\label{tab:rq1-timing}
\small
\begin{tabular}{l l r r}
\toprule
\textbf{Case} & \textbf{Error Pattern} & \textbf{Mean (s)} & \textbf{Range (s)} \\
\midrule
div\_in\_path\_WANGMI             & Div In Path              & 14.94 & 12.0--16.8 \\
web3bugs\_45\_H\_01               & Erroneous Acct.          & 13.55 & 9.4--18.4  \\
web3bugs\_58\_H\_02               & Erroneous Acct.          &  9.67 & 7.0--13.8  \\
indivisible\_amount\_BoostToken   & Minor Amt.\ Retention    &  5.67 & 5.0--6.1   \\
operator\_order\_issue\_BoostToken& Operator Order Issue      &  5.43 & 4.9--5.8   \\
web3bugs\_62\_H\_08               & Inconsist.\ State Upd.   &  4.76 & 4.4--5.1   \\
web3bugs\_70\_H\_10               & Inconsist.\ State Upd.   &  3.32 & 3.0--3.6   \\
web3bugs\_51\_H\_02               & Erroneous Acct.          &  3.18 & 2.8--4.0   \\
web3bugs\_56\_H\_02               & Erroneous Acct.          &  2.37 & 1.5--4.0   \\
profit\_opportunity\_HIT          & Exchange Problem         &  2.03 & 1.6--2.5   \\
web3bugs\_5\_H\_12                & Erroneous Acct.          &  1.75 & 0.9--3.4   \\
web3bugs\_47\_H\_02               & Erroneous Acct.          &  1.42 & 1.3--1.5   \\
web3bugs\_5\_H\_08                & Erroneous Acct.          &  1.10 & 1.0--1.3   \\
web3bugs\_60\_H\_01               & Erroneous Acct.          &  1.06 & 1.0--1.1   \\
web3bugs\_101\_H\_01              & Erroneous Acct.          &  0.87 & 0.8--0.9   \\
exchange\_problem\_Nokon          & Exchange Problem         &  0.67 & 0.6--0.8   \\
greedy\_contract\_SwordCrowdsale  & Greedy Contract          &  0.66 & 0.6--0.7   \\
web3bugs\_78\_H\_02               & Erroneous Acct.          &  0.58 & 0.6--0.6   \\
web3bugs\_5\_H\_07                & Erroneous Acct.          &  0.44 & 0.3--0.6   \\
web3bugs\_77\_H\_01               & Erroneous Acct.          &  0.33 & 0.3--0.4   \\
\midrule
\textbf{Mean of means} & & \textbf{3.69} & --- \\
\bottomrule
\end{tabular}
\end{table}

\begin{figure}[t]
\centering
\includegraphics[width=\linewidth]{figure/rq1_scatter.pdf}
\caption{Scatter plots of per-case analysis time against six complexity metrics (20 mitigated cases, 3-run mean). Dashed lines show linear regression; $r$ is the Pearson correlation coefficient.}
\label{fig:rq1-scatter}
\end{figure}

\noindent\textbf{Setup.}
For each of the 75 effective benchmark instances (§\ref{sec:setup}), we (1) read the accompanying bug documentation and identify the buggy line(s), (2) write @During and/or @Post intent annotations that capture the expected numeric behavior described therein, (3) supply the debug annotations (@StateVar, @LocalVar, @GlobalVar, and @IReturn) that the function needs to be analyzable under representative inputs, and (4) run IntentChecker on the resulting annotated source.
We record the judgment produced by the validation engine---Satisfied, Warning, or Violated---together with the risk score and violated region.
For instances where no annotation drives the validation engine to Violated on the buggy code, we further classify the structural reason according to the limitation taxonomy in RQ2 (§\ref{sec:rq2}).

\smallskip
\noindent\textbf{Outcomes.}
Of the 75 effective instances, IntentChecker produces a Violated judgment on \textbf{20 cases} when supplied with annotations derived from the bug documentation.
The remaining 55 instances are classified as not\_detectable: no annotation can drive the validation engine to Violated on the buggy code, given the structure of its analysis or its intent model.
Table~\ref{tab:rq1-overview} summarizes the full outcome distribution; the structural reasons behind the 55 unmitigated cases are analyzed in RQ2 (§\ref{sec:rq2}).

\smallskip
\noindent\textbf{Annotation model usage on the 20 mitigated cases.}
The upper half of Table~\ref{tab:rq1-overview} reports which intent annotations were used to surface the bug in each mitigated case.
Of the 20 cases, 12 use a @During annotation (mid-execution checks) and 8 use a @Post annotation (post-execution checks).
The single most common form is the simple relational comparison (intentValue relOp intentValue), used in 9 cases across both @During and @Post.
The two specialized @During forms---Before/After comparisons against an injection point, and func.arg[N] checks on argument values---are each used 2--3 times, and the specialized @Post forms (Entry/Exit, returnExpression, changed(\,)) account for the remainder.
This distribution suggests that the bulk of mitigation work in our benchmark is carried by the basic relational form, with the specialized annotations contributing on patterns that the basic form cannot capture (e.g., monotonicity over function execution, argument-level constraints, post-execution change detection).

\smallskip
\noindent\textbf{Pattern coverage of the 20 mitigated cases.}
As shown in the Error Pattern column of Table~\ref{tab:rq1-timing}, the 20 mitigated instances cover 7 of the 8 single-transaction error patterns enumerated in §\ref{sec:background}.1: Div In Path, Operator Order Issue, Exchange Problem, Minor Amount Retention, Erroneous Accounting (with both Wrong Interest Rate Order and Wrong Checkpoint Order sub-patterns appearing), Inconsistent State Updates, and Greedy Contract.
Precision Loss Trend is present in the benchmark (2 cases) but both instances are not mitigated.
The coverage across 7 of 8 patterns indicates that mitigability depends on whether the error is expressible as an intent violation, not on the error's classification---suggesting that intent-based validation is pattern-agnostic in scope.

\smallskip
\noindent\textbf{Analysis time and what drives it.}
We measured per-case end-to-end analysis time over three independent runs of each mitigated case on a single thread, excluding the one-time dependency pre-analysis phase (inherited contracts and libraries are pre-analyzed and their control-flow graphs are stored as serialized objects, so they need not be rebuilt across editing sessions).
The mean per-case time is \textbf{3.69\,s}, the minimum is 0.33\,s (web3bugs\_77\_H\_01), and the maximum is 14.94\,s (div\_in\_path\_WANGMI); 15 of the 20 cases finish in under 5\,s, and all but two finish in under 10\,s.
The full per-case timings appear in Table~\ref{tab:rq1-timing}.

To understand what drives the variation across cases, we collected five per-case complexity metrics and computed the Pearson correlation coefficient $r$~\citep{pearson_r} of each with the 3-run mean analysis time $\bar{y}_i$.
For two variables $X$ (metric) and $Y$ (analysis time) with $n$ paired observations $\{(x_i,y_i)\}_{i=1}^{n}$,
\[
  r \;=\; \frac{\sum_{i}(x_i-\bar{x})(y_i-\bar{y})}
               {\sqrt{\sum_{i}(x_i-\bar{x})^{2}}\;\cdot\;\sqrt{\sum_{i}(y_i-\bar{y})^{2}}}\,,
  \qquad -1 \le r \le +1\,,
\]
where $r=+1$ indicates a perfect positive linear relationship and $r=0$ indicates no linear association.

Figure~\ref{fig:rq1-scatter} plots each metric against analysis time.
Four of the five metrics exhibit a strong positive correlation ($r>0.8$):
\begin{itemize}
  \item \textbf{Source lines} ($r=0.87$) is the strongest predictor, as it directly reflects the amount of code the abstract interpreter must traverse.
  \item \textbf{Control flow} ($r=0.84$), the aggregate count of loops, branches, and guards (require/assert), increases the number of join operations and fixpoint iterations the abstract interpreter must perform.
  \item \textbf{Internal calls} ($r=0.80$) captures the depth of the call chain within the contraction. Each internal call requires the abstract interpreter to enter a new function scope, propagate abstract states across the call boundary, and merge the return state back into the caller.
  \item \textbf{Debug annotations} ($r=0.80$) adds tracked variables to the abstract state, increasing per-step cost proportionally.
\end{itemize}
The consistently strong correlation across these four metrics ($r \ge 0.80$, $n=20$) suggests that analysis time scales approximately linearly with the size and complexity of the analysis input within the range observed in our benchmark.
\begin{itemize}
  \item \textbf{External calls} ($r=0.69$) shows a comparatively weaker correlation. Dependencies are excluded from per-case time as noted above, and interface function calls lack an implementation body, so the analyzer resolves them to $\top$ or a user-supplied @IReturn value without traversal.
\end{itemize}

Taken together, RQ1 shows that intent-based validation can surface actionable diagnostic signals for 20 of 75 benchmark instances across 7 of 8 error patterns, with per-case analysis completing in under 15\,s (mean 3.69\,s) and scaling proportionally with code complexity.
These results suggest that IntentChecker is practical as a development-time validation tool for numeric logic errors in smart contracts.

\subsection{RQ2: Unmitigated Case Analysis}\label{sec:rq2}

\noindent\textbf{Setup.}
RQ2 classifies the 55 instances for which IntentChecker cannot produce a Violated judgment on the buggy code into analysis-engine limitations (L1--L3) and annotation-grammar limitations (L4--L5), and further examines the expressibility of the annotation grammar within L4--L5 (Figure~\ref{fig:rq1-limitations}).

\begin{figure}[t]
\centering
\begin{tikzpicture}[font=\small]
  % Constants
  \def\barH{0.45}
  \def\barGap{0.18}
  \def\xUnit{0.55}  % cm per case (widened)
  \def\maxCount{10}

  % Group A: analysis-engine (L1-L3)
  \def\colA{black!55}
  % Group B: annotation-grammar (L4-L5)
  \def\colB{black!25}

  % Bars: (index/label/count/color/type name)
  \foreach \i/\lab/\val/\col/\desc in {%
    0/L1a/8/\colA/Loop widening,
    1/L1b/2/\colA/Loop body granularity,
    2/L2a/3/\colA/Interface call return $\top$,
    3/L2b/1/\colA/External call unknown state,
    4/L3/7/\colA/Unsupported construct,
    5/L4a/10/\colB/Inexpressible expected value,
    6/L4b/8/\colB/No target storage variable,
    7/L4c/1/\colB/Magnitude-only difference,
    8/L4d/1/\colB/Invariant masked,
    9/L5a/7/\colB/Bug-aware: missing code,
    10/L5b/7/\colB/Bug-aware: wrong code%
  } {%
    \pgfmathsetmacro{\y}{-\i*(\barH+\barGap)}
    \pgfmathsetmacro{\w}{\val*\xUnit}
    \node[anchor=east] at (0,\y) {\textbf{\lab}};
    \fill[\col] (0.15,\y-\barH/2) rectangle (0.15+\w,\y+\barH/2);
    \node[anchor=west] at (0.18+\w,\y) {\val\;\scriptsize\desc};
  }

  % Group brackets
  \draw[thick] (-1.05,0.32) -- (-1.15,0.32) -- (-1.15,-2.06) -- (-1.05,-2.06);
  \node[rotate=90, anchor=south] at (-1.25,-0.87) {\scriptsize Analysis-engine (21)};
  \draw[thick] (-1.05,-2.46) -- (-1.15,-2.46) -- (-1.15,-6.84) -- (-1.05,-6.84);
  \node[rotate=90, anchor=south] at (-1.25,-4.65) {\scriptsize Annotation-grammar (34)};

  % X-axis
  \draw[->] (0.15,-7.25) -- (0.15+11*\xUnit,-7.25);
  \foreach \x in {0,2,4,6,8,10} {
    \draw (0.15+\x*\xUnit,-7.20) -- (0.15+\x*\xUnit,-7.30);
    \node[anchor=north] at (0.15+\x*\xUnit,-7.30) {\scriptsize \x};
  }
  \node[anchor=north] at (0.15+5*\xUnit,-7.65) {\scriptsize \# of cases};
\end{tikzpicture}
\caption{Limitation breakdown for the 55 not\_detectable instances: analysis-engine (L1--L3, 21) vs.\ annotation-grammar (L4--L5, 34).}
\label{fig:rq1-limitations}
\end{figure}

\smallskip
\noindent\textbf{Limitation type definitions.}
\begin{itemize}
  \item \textbf{Analysis-engine limitations (L1--L3, 21 cases).} A bug-relevant value is abstracted to $\top$ during analysis, so the interpreter cannot distinguish the buggy and correct executions, and no developer annotation can recover the lost precision.
  \begin{itemize}
    \item \textbf{L1a (Loop widening, 8).} A variable that the bug check depends on is widened to $\top$ inside a preceding loop, and the $\top$ propagates past the loop exit to the check point. The widened interval admits both the buggy and correct outcomes, so the annotation evaluates to Warning rather than Violated---an ambiguous judgment that offers no actionable signal distinguishing the buggy code from the corrected one.
    \item \textbf{L1b (Loop body granularity, 2).} The bug manifests only at a specific iteration of a loop, but the analyzer's state inside the loop body is a single summary joined over all iterations, so an @During annotation placed there cannot single out the faulty iteration.
    \item \textbf{L2a (Interface call return $\top$, 3).} The bug check depends on a value derived from a state-modifying interface call. The implementation is unavailable and @IReturn applies only to view/pure calls (§\ref{sec:prelim:extension}(b)), so the return collapses to $\top$ and propagates to the check point.
    \item \textbf{L2b (External call, unknown external state, 1).} The bug check depends on the return of a cross-contract call whose implementation is available through dependency pre-analysis. An L1, L2a, or L3 pattern within the callee's own body causes $\top$ to propagate to its return value, which then reaches the check point.
    \item \textbf{L3 (Unsupported construct, 7).} The bug check depends on a value that flows through an unsupported construct---abi.decode, opaque hash builtins, or non-arithmetic Yul---which the analyzer treats as $\top$, and the $\top$ propagates to the check point.
  \end{itemize}
  \noindent These 21 cases are residual failures that remain after the analyzer's precision mechanisms---annotation-guided adaptive widening inherited from SolQDebug~\citep{solqdebug}, debug-annotation-supplied concrete values, and dependency pre-analysis for cross-contract resolution---have been applied. The presence of a loop, an external call, or a low-level construct does not by itself block analysis. L1--L3 arise only when the bug-relevant value flows through the specific failure modes enumerated above.
  \item \textbf{Annotation-grammar limitations (L4--L5, 34 cases).} The abstract interpreter would distinguish the buggy and correct executions if the developer could express the right intent, but the current grammar either has no such expression (L4) or has one whose construction presupposes already knowing the bug (L5).
  \begin{itemize}
    \item \textbf{L4a (Inexpressible expected value, 10).} The correct expected value depends on a new intermediate computation, an external contract's state, or a function call that does not appear in the program; no intentValue expression can be constructed.
    \item \textbf{L4b (No target storage variable, 8).} The buggy function (a view function, a pure function, a thin wrapper, or a library helper) does not modify any storage variable in the target contract, so no @Post annotation has anything to attach to.
    \item \textbf{L4c (Magnitude-only difference, 1).} Buggy and correct code change the same state variable in the same direction; only the magnitude differs. changed(\,) and Entry op Exit both succeed in both versions.
    \item \textbf{L4d (Invariant masked, 1).} A different statement in the same function already modifies the target state variable, so changed(\,) is satisfied in both versions. The actual invariant is a multi-variable relation that PostEntryExit cannot express.
    \item \textbf{L5a (Bug-awareness-required: missing code, 7).} A required statement is omitted; the post-condition that would catch this is expressible, but knowing what is missing already presupposes knowing the bug.
    \item \textbf{L5b (Bug-awareness-required: wrong code, 7).} A required statement uses the wrong identifier, operator, or struct field; the corrective annotation is expressible, but writing it requires knowledge of the correct form.
  \end{itemize}
\end{itemize}

\noindent\textbf{Expressibility within L4--L5.}
Within the 34 annotation-grammar limitation cases, we further ask: assuming full bug awareness, could a developer construct an annotation that the validation engine would surface as Violated or Warning?
L4's 20 cases are inexpressible by definition---no annotation under the current grammar can yield a non-Satisfied judgment.
All 14 L5 cases, however, are expressible, splitting evenly into two tiers:

\begin{itemize}
  \item \textbf{Expressible-Strong (``Direct'', 7).} The annotation expresses the correct value or relation directly. The buggy version produces Violated; the corrected version produces Satisfied. The annotation is essentially a fix-equivalent specification.
  \item \textbf{Expressible-Weak (``Indirect'', 7).} The annotation captures only that some state variable changed---or did not change---in a particular direction (changed(var, true/false) or var(Entry relOp Exit)). The buggy version produces Violated, but the annotation does not encode the correct value, only its presence or polarity---an awareness signal rather than a fix-equivalent specification.
\end{itemize}

\smallskip
\noindent\textbf{Representative examples.}
We illustrate one case from each tier to show the qualitative difference in signal.

\noindent Inexpressible (L4a) --- web3bugs\_51\_H\_04 (SwapUtils.getYC). The bug is that when a swap crosses a target price, the StableSwap formula switches its amplifier coefficient mid-trade; the correct behavior is to split the swap at the crossing point and apply each amplifier to its sub-swap. Constructing this requires computing a split point $dx_1$ that does not appear anywhere in the program and that is not expressible as a closed-form arithmetic combination of in-scope variables (it is the solution to a Newton-method equation). No intentValue expression can capture the correct getY return, so the annotation grammar's intentValue == arithExpr form is unusable here.

\noindent Expressible-Strong (L5b) --- web3bugs\_113\_H\_05 (NFTPairWithOracle.\_lend). The buggy line is require(params.ltvBPS $\geq$ accepted.ltvBPS, ...); the correct comparison is $\leq$ (a lender wants the borrower's LTV to be at most their accepted ceiling). Under bug awareness, the annotation is direct:
\begin{verbatim}
// @During params.ltvBPS <= accepted.ltvBPS
\end{verbatim}
The buggy code admits inputs with params.ltvBPS $>$ accepted.ltvBPS, which violate the annotation cleanly. The annotation is essentially the corrective require.

\noindent Expressible-Weak (L5a) --- web3bugs\_192\_H\_01 (Lock.extendLock). The function performs transferFrom for an additional \_amount but forgets to update the internal accounting variable totalLocked[\_asset]. Under bug awareness, the developer can write
\begin{verbatim}
// @Post totalLocked[_asset](Entry < Exit)
\end{verbatim}
which fires Violated on the buggy version (the variable does not change, so Entry == Exit) and is satisfied on a corrected version that adds totalLocked[\_asset] += \_amount. The annotation tells the developer that the variable was supposed to grow, but it does not tell them by how much---the precise increment \_amount is not present in the annotation. This is an awareness signal, not a fix-equivalent specification, and our grammar does not currently support intentValue expressions over the Entry value, which would be needed to encode the increment exactly.

\smallskip
\noindent\textbf{Implications.}
Three findings stand out.
First, L4's 20 cases are a hard ceiling that no amount of developer awareness can lift, whereas L5's 14 cases are a soft ceiling that the intent model already covers, contingent on whether the developer is aware.
Second, within L5, the even split between Strong and Weak suggests that for half of the awareness-conditional cases, the developer's annotation would express the actual fix, and for the other half it would express a coarser ``something must change here'' signal that nevertheless suffices to surface the bug.
Third, the L4 cases reveal a fundamental boundary of abstraction level: IntentChecker's intent model operates at the value level---it can express whether a variable's value satisfies a numeric relationship---but it cannot express whether the right computation was performed in the first place.
In L4a, the correct expected value requires a computation that does not exist in the program; in L4b, the function modifies no state variable, so there is no value to assert on.
The developer's true intent in these cases is about what the code should \textit{do} (``split the swap at the crossing point'', ``update totalLocked''), not about what value a variable should hold.
Capturing such logic-level intent would require a fundamentally different formalism, such as temporal properties over execution traces or structural specifications of required code patterns.
IntentChecker is effective when the bug manifests as a wrong value; complementary tools are needed when the bug manifests as a missing or misplaced operation.
Future grammar extensions---specifically, (a) arithmetic expressions over Entry values in PostEntryExit, which would absorb several Weak cases, and (b) a returnExpression clause for view functions, which would absorb several L4b cases---are concrete directions discussed in §\ref{sec:disc:tech}.

\subsection{RQ3: Comparison with Existing Detection Tools}\label{sec:rq3}

\begin{table}[t]
\centering
\caption{Signal characterization of the four compared tools.}
\label{tab:rq3-signal}
\small
\begin{tabular}{@{}lllll@{}}
\toprule
 & IntentChecker & ScType & GPTScan & NumScout \\
\midrule
Approach & Abstract interp. & Type inference & LLM + static & Symbolic exec. \\
Detection model & Specification-driven & Type-driven & Pattern-driven & Pattern-driven \\
Annotation & During/Post conds. & Financial types & None & None \\
Target class & Developer-specified & Type inconsistency & DeFi logic (10 rules) & Numeric (5 patterns) \\
\bottomrule
\end{tabular}
\end{table}

\begin{table}[t]
\centering
\caption{Mitigation results on the 20 RQ1-mitigated cases. \ding{51}\,=\,mitigated (tool produced an actionable signal for the target bug), \ding{55}\,=\,not mitigated (tool was applicable but produced no signal), \textcolor{gray}{N/A}\,=\,out of scope (no applicable rule or type file, or execution infeasible). GPTScan entries report results as published in Sun~et~al.~\cite{gptscan} rather than reruns; only 45\_H\_01 is listed there as a true positive (MSA-002, Wrong Order for Interest on UToken.borrow), and the remaining 19 cases either lie in contests outside that paper's evaluation set (51, 58, 60, 62, 101) or involve bug classes (wrong-formula, wrong-scale, wrong-decimals, missing-update) that do not match any of GPTScan's ten rules. NumScout entries marked $^{\dagger}$ were evaluated on a semantically-equivalent variant in which external library functions were rewritten as internal to bypass NumScout's library-placeholder limitation; see §\ref{sec:rq3} for details.}
\label{tab:rq3-mitigation}
\scriptsize
\begin{tabular}{@{}llcccc@{}}
\toprule
Case & Contract.Function & IC & ScType & GPTScan & NumScout \\
\midrule
WANGMI & WANGMI.\_transfer & \ding{51} & \textcolor{gray}{N/A} & \textcolor{gray}{N/A} & \ding{51} \\
Nokon & Nokon.buy & \ding{51} & \textcolor{gray}{N/A} & \textcolor{gray}{N/A} & \ding{51} \\
SwordCrowdsale & SwordCrowdsale.refundMoney & \ding{51} & \textcolor{gray}{N/A} & \textcolor{gray}{N/A} & \ding{55} \\
BoostToken (op) & BoostToken.sendETHToTeam & \ding{51} & \textcolor{gray}{N/A} & \textcolor{gray}{N/A} & \ding{51} \\
BoostToken (ind) & BoostToken.sendETHToTeam & \ding{51} & \textcolor{gray}{N/A} & \textcolor{gray}{N/A} & \ding{51} \\
HIT & HIT.getTokens & \ding{51} & \textcolor{gray}{N/A} & \textcolor{gray}{N/A} & \ding{51} \\
5\_H\_07 & Utils.calcAsymmetricShare & \ding{51} & \ding{51} & \textcolor{gray}{N/A} & \ding{51} \\
5\_H\_08 & Utils.calcLiquidityUnits & \ding{51} & \ding{51} & \textcolor{gray}{N/A} & \ding{51} \\
5\_H\_12 & Pools.getAddedAmount & \ding{51} & \textcolor{gray}{N/A} & \textcolor{gray}{N/A} & \ding{55} \\
45\_H\_01 & UToken.borrow & \ding{51} & \textcolor{gray}{N/A} & \ding{51} & \ding{51} \\
47\_H\_02 & WrappedIbbtcEth.transferFrom & \ding{51} & \ding{51} & \textcolor{gray}{N/A} & \ding{55} \\
51\_H\_02 & SwapUtils.rampTargetPrice & \ding{51} & \textcolor{gray}{N/A} & \textcolor{gray}{N/A} & \ding{55}$^{\dagger}$ \\
56\_H\_02 & CDP.update & \ding{51} & \ding{55} & \textcolor{gray}{N/A} & \ding{55} \\
58\_H\_02 & LpIssuer.\_chargeFees & \ding{51} & \textcolor{gray}{N/A} & \textcolor{gray}{N/A} & \ding{55} \\
60\_H\_01 & OptimisticLedgerLib.settleAccount & \ding{51} & \ding{51} & \textcolor{gray}{N/A} & \ding{55} \\
62\_H\_08 & Stream.updateStreamInternal & \ding{51} & \textcolor{gray}{N/A} & \textcolor{gray}{N/A} & \ding{55} \\
70\_H\_10 & LiquidityBasedTWAP.syncVaderPrice & \ding{51} & \ding{55} & \textcolor{gray}{N/A} & \ding{55} \\
77\_H\_01 & MathLib.calcLiqTokenQty & \ding{51} & \textcolor{gray}{N/A} & \textcolor{gray}{N/A} & \ding{51}$^{\dagger}$ \\
78\_H\_02 & RebaseProxy.mint & \ding{51} & \textcolor{gray}{N/A} & \textcolor{gray}{N/A} & \ding{55} \\
101\_H\_01 & LenderPool.\_calcPrincipalWithdrawable & \ding{51} & \ding{55} & \textcolor{gray}{N/A} & \ding{55} \\
\bottomrule
\end{tabular}
\end{table}

We compare IntentChecker against three state-of-the-art tools on the 20 mitigated cases from RQ1:
\textbf{ScType}~(ICSE~'24), a static type-checker that assigns financial types (token unit, scaling factor) to variables and detects type inconsistencies;
\textbf{GPTScan}~(ICSE~'24), an LLM-based multi-step pipeline that matches functions against ten DeFi-specific vulnerability rules (price manipulation, slippage, etc.); and
\textbf{NumScout}, a symbolic-execution engine that searches for five predefined numeric defect patterns (division-in-path, operator order, exchange problem, precision loss, minor amount retention).
Each tool was executed three times on the same machine; ScType and NumScout produce deterministic output, while GPTScan's LLM-based pipeline may vary across runs.

\smallskip
\noindent\textbf{NumScout execution adjustments.}
To run NumScout over our Web3Bugs subset we applied two classes of adjustments, both external to the bugs under study.
First, a source-level bug in NumScout's semantic analyzer (feature\_detector/semantic\_analysis.py, in the safeTransfer handler) raised UnboundLocalError on every contract reaching the relevant branch; we applied a minimal one-line fix (renaming an undefined local to the intended stack variable) so the tool could complete analysis.
Second, NumScout consumes EVM bytecode and fails with ``incomplete push instruction'' errors when the bytecode contains unresolved library placeholders (\_\_\$\ldots\$\_\_), which are emitted when a contract calls a public/external library function that would normally be linked at deploy time.
Two of our cases (51\_H\_02, Swap.rampTargetPrice; 77\_H\_01, Exchange.calculateLiquidityTokenQtyForSingleAssetEntry) use external libraries in this way.
For these two cases only, we produced a second copy of the target file in which every public/external library function was rewritten as internal, causing the compiler to inline the library body into the calling contract's bytecode rather than emit a placeholder.
This transformation is semantically equivalent at the operation level---the arithmetic instructions NumScout inspects (MUL, DIV, SUB, etc.), the storage accesses, and the control flow are preserved---so the defect patterns NumScout targets are unchanged; only the call-site mechanism (DELEGATECALL vs.\ inlined code) differs.
The modified files are kept separate from the originals used in the ScType and GPTScan experiments, and both modifications are restricted to the two cases noted above.

\smallskip
\noindent\textbf{Detection model.}
A key distinction among these tools is how the set of detectable bugs is defined.
GPTScan and NumScout operate on a pattern-driven model: each tool ships with a fixed catalog of vulnerability patterns (10 rules and 5 patterns, respectively), and only bugs matching those patterns can be detected.
When a new class of numeric error emerges, the catalog must be updated by the tool developers.
ScType is type-driven: it checks whether arithmetic operations preserve financial type consistency, which is more general than a fixed pattern catalog but still limited to errors expressible as type mismatches.
IntentChecker is specification-driven: developers write @During/@Post conditions describing intended behavior, and the tool verifies the implementation against those conditions via abstract interpretation.
This allows IntentChecker to target any numeric property expressible within its annotation grammar---including novel error patterns not yet cataloged by existing tools---but requires developer participation.
The grammar's expressibility ceiling was quantified in RQ2.

\begin{figure}[t]
    \centering
    \includegraphics[width=\linewidth]{figure/rq3_time_comparison.pdf}
    \caption{Analysis time comparison across four tools (log scale). Left: box plot of per-case times over three runs. Right: mean analysis time. IntentChecker and ScType complete in seconds; GPTScan and NumScout each require minutes per case, with NumScout exceeding GPTScan on cases that exercise its symbolic-execution backend.}
    \label{fig:rq3-time}
\end{figure}

\smallskip
\noindent\textbf{Finding 1: coverage landscape.}
Table~\ref{tab:rq3-mitigation} shows the per-case mitigation results.
For GPTScan we report the per-contest results published in Sun~et~al.~\cite{gptscan} rather than reruns, because GPTScan's source invokes OpenAI's \texttt{gpt-3.5-turbo} alias whose underlying snapshot has been redirected multiple times since 2023 (\S\ref{sec:threats}).
Seven of the twelve Web3Bugs contests in our benchmark (5, 45, 47, 56, 70, 77, 78) fall within GPTScan's evaluation set; the other five (51, 58, 60, 62, 101) are not evaluated in that paper.
Within the overlap, only 45\_H\_01 is listed as a true positive---a Wrong Order for Interest hit on \texttt{UToken.borrow}---and the remaining overlapping bugs are not categorized under any of GPTScan's ten rules in the published breakdown.
The nineteen non-matched cases involve wrong-formula, wrong-scale, wrong-decimals, and missing-update errors that do not correspond to GPTScan's patterns (price manipulation, slippage limit, first deposit, wrong-order interest/checkpoint, etc.).
This is not a deficiency of GPTScan---the tool was designed for a different threat model---but it demonstrates that a coverage gap exists: the bug class targeted by IntentChecker is largely outside the reach of the most recent LLM-based smart-contract analysis tool.

NumScout mitigated the target bug in 9 of 20 cases.
The nine mitigated cases concentrate on NumScout-origin samples (WANGMI, Nokon, BoostToken, HIT) and on Vader Protocol's Utils.sol, whose division-heavy arithmetic matches NumScout's division-in-path and operator-order patterns; Web3Bugs cases whose bugs lie outside these five patterns (e.g., CDP.update, Stream.updateStreamInternal, LiquidityBasedTWAP.syncVaderPrice, RebaseProxy.mint) are not mitigated despite NumScout completing analysis.
For 51\_H\_02 and 77\_H\_01, the reported NumScout entries are obtained on the library-inlined variant described earlier, and are marked with a dagger in Table~\ref{tab:rq3-mitigation}.

\smallskip
\noindent\textbf{Finding 2: complementarity with ScType.}
ScType provides the most direct comparison, as both tools target accounting errors and both require manual annotations (financial types for ScType, during/post conditions for IntentChecker).
Among the seven overlapping cases, ScType mitigated four: calcAsymmetricShare, calcLiquidityUnits (both in Vader Protocol's Utils.sol), transferFrom (WrappedIbbtcEth), and settleAccount (OptimisticLedger).
IntentChecker mitigated all seven.
The three cases that ScType could not mitigate (CDP.update, syncVaderPrice, \_calcPrincipalWithdrawable) are not type-inconsistency errors but rather value-level errors: the arithmetic uses correctly-typed operands but produces a numerically wrong result.
ScType's type-driven model cannot express the property ``this subtraction should yield the correct fee amount''; IntentChecker's specification-driven model can.
Conversely, ScType's type inference operates without requiring the developer to know the expected output value, which makes it applicable even when the developer does not have a specific numeric property in mind.
The two approaches are thus complementary: ScType catches type-level mismatches automatically, while IntentChecker catches value-level deviations when the developer can articulate the intended behavior.

\smallskip
\noindent\textbf{Finding 3: analysis latency.}
Figure~\ref{fig:rq3-time} reports per-case analysis times averaged over three runs.
IntentChecker (mean 4.5\,s) and ScType (mean 13.9\,s) are both local, deterministic analyses that complete in seconds.
GPTScan (mean 267\,s) is approximately 59$\times$ slower due to its multi-step LLM pipeline, and NumScout (mean 772\,s, median 517\,s over 20 cases) is slower still, with several Web3Bugs cases running to the 1800\,s global timeout set by NumScout's symbolic-execution backend.
GPTScan's per-case latency is dominated by sequential OpenAI API round-trips over its two-step LLM filter, and its pipeline structure has not changed since 2023; the order-of-magnitude comparison (seconds vs.\ minutes per case) is therefore insensitive to the model-snapshot changes that affect mitigation outputs (§\ref{sec:threats}).
The sub-10-second latency of IntentChecker makes it suitable for integration into CI/CD pipelines or interactive editor feedback loops, where developers expect near-instant responses.



\smallskip
\noindent\textbf{Summary.}
Each tool occupies a distinct position in the detection landscape:
pattern-driven tools (GPTScan, NumScout) offer fully automatic analysis but can only mitigate bugs within their predefined vulnerability catalogs;
type-driven analysis (ScType) generalizes beyond fixed patterns but is limited to errors expressible as type mismatches;
and specification-driven analysis (IntentChecker) can mitigate any property within its annotation grammar but requires developer participation.
The 20-case evaluation demonstrates that numeric logic errors---a class responsible for some of the highest-impact DeFi exploits---fall entirely outside the scope of pattern-driven tools and partially outside the reach of type-driven analysis, confirming the practical need for a specification-driven complement.

\subsection{Threats to Validity}\label{sec:threats}

\noindent\textbf{Internal validity.}
The intent and debug annotations used in RQ1 were authored by the researchers based on the audit report descriptions, not by the original contract developers.
A different annotator could plausibly write a slightly different annotation for the same bug, and a sufficiently weak annotation could fail to surface a bug that a stronger annotation would catch.
We mitigate this by deriving each annotation directly from the audit report's stated root cause and recording the exact annotation alongside each case in the supplementary material, allowing independent re-evaluation.

\smallskip
\noindent\textbf{External validity.}
Our 89-case benchmark is drawn from Web3Bugs (Code4rena audit findings) and NumScout (BNB Smart Chain), which together cover a substantial portion of publicly known numeric logic errors but are not exhaustive.
Findings may not generalize to numeric logic errors that arise in domains underrepresented by these two sources, such as zero-knowledge circuits or non-DeFi applications.
Additionally, the correlation analysis in RQ1 is based on $n=20$ mitigated cases, which limits statistical power; the reported Pearson coefficients should be interpreted as indicative trends rather than causal claims.

\smallskip
\noindent\textbf{Construct validity.}
RQ1's notion of ``mitigation'' relies on the validation engine producing a non-Satisfied judgment (typically Violated) for the annotated bug, rather than on a downstream measurement of developer behavior.
A non-Satisfied judgment on its own does not guarantee that a developer reading it would correctly diagnose and fix the bug; assessing that downstream effect is an important direction for future work.
In RQ3, GPTScan results are taken from the per-contest breakdown published in Sun~et~al.~\cite{gptscan} rather than from reruns in our environment.
GPTScan's implementation invokes the model alias \texttt{gpt-3.5-turbo} (and \texttt{gpt-4}) whose underlying snapshots have been redirected by the provider multiple times since the original evaluation in 2023; our April~2026 reruns of the same code on the same inputs produced rule activations at different functions than the published results, confirming that model-snapshot drift alters the tool's output even when code and inputs are held constant.
We therefore report the authors' published numbers as the fairest reflection of GPTScan's intended performance.
The per-case latency in Figure~\ref{fig:rq3-time} still uses our reruns because per-case timings were not published in the original paper; this affects only the absolute latency comparison, not the per-case mitigation entries.
Cases marked N/A for GPTScan in Table~\ref{tab:rq3-mitigation} fall into two categories: (i) contests outside the evaluation set of Sun~et~al.~(51, 58, 60, 62, 101); and (ii) bugs whose class is not among GPTScan's ten rules (wrong-formula, wrong-scale, wrong-decimals, missing-update).
The ScType comparison is limited to seven overlapping cases for which type annotation files exist; results on the remaining thirteen cases may differ.
Two NumScout cases (51\_H\_02, 77\_H\_01) were evaluated on a variant of the target file in which public/external library functions were rewritten as internal to bypass NumScout's inability to handle unresolved library placeholders in compiled bytecode (§\ref{sec:rq3}).
Because this transformation only alters the call-site mechanism (inlined library code vs.\ DELEGATECALL) and preserves the arithmetic instructions, storage accesses, and control flow that NumScout's defect patterns inspect, we regard the comparison as fair for the specific pattern-level signal NumScout produces; we nevertheless note the modification explicitly so that readers interpreting Table~\ref{tab:rq3-mitigation} can account for it.

\section{Discussion}\label{sec:discussion}

\subsection{Annotation Authoring Burden and Future Automation}\label{sec:disc:burden}

IntentChecker's mitigation results in §\ref{sec:rq1} are conditional on the developer supplying intent and debug annotations.
This is by design---the entire approach is predicated on the developer making their intent explicit---but it nevertheless constitutes the tool's most significant practical limitation.
Two distinct kinds of effort are involved.

\smallskip
\noindent\textbf{Effort 1: writing intent annotations for unfamiliar code.}
The intent model assumes that the developer can articulate, in the grammar of @During/@Post annotations, what the function is supposed to do.
For code the developer wrote themselves a few minutes ago, this is unproblematic: the intent is still in the developer's head.
For code the developer is reviewing, modifying, or debugging without prior context, the intent must first be reconstructed from the surrounding code, comments, and external specifications, and only then can it be encoded as an annotation.
This reconstruction step is not free, and it is exactly the step that the L5 (bug-awareness-required) limitation in §\ref{sec:rq2} measures: when the correct intent is non-obvious, an annotation that captures it cannot be written without already knowing the bug.

\smallskip
\noindent\textbf{Effort 2: supplying debug annotations for realistic execution paths.}
Even for code whose intent is clear, IntentChecker can only analyze the path that the developer's debug annotations actually exercise.
@StateVar, @LocalVar, and @IReturn require the developer to choose representative ranges; choosing ranges that miss the buggy region produces a false negative, while choosing ranges that are too wide invites widening and produces uninformative Warning judgments.
This is the standard cost of any annotation-based analysis and is well understood, but it remains a non-trivial barrier to adoption.

\smallskip
\noindent\textbf{The circularity objection.}
A natural question is whether writing a correct intent annotation already presupposes knowing the bug---in which case the tool would only confirm what the developer already knows.
The concern is legitimate, and our RQ2 analysis quantifies it precisely: 14 of the 34 intent-model limitation cases (L5) require bug awareness before the annotation can be written, confirming that for those cases, independent discovery through annotation alone is not realistic.
However, this circularity is not unique to IntentChecker; it is inherent in any specification-based approach, from formal pre/post-conditions to Scribble annotations to test oracles.
What differs across tools is how much of the specification burden falls on the developer.
For the 20 mitigated cases in RQ1, the annotations were derivable from the developer's own stated intent (e.g., ``this fee should equal $x$'', ``this balance should decrease after withdrawal'') without requiring knowledge of the specific bug.
The L5 cases, by contrast, identify the boundary where intent-driven annotation reaches its limit and external assistance---such as automated annotation suggestion---becomes necessary.
This observation directly motivates the LLM-assisted workflow described below: rather than expecting the developer to write all annotations from scratch, a suggestion layer can propose candidate annotations that the developer accepts, rejects, or refines, shifting the task from generation to validation and partially breaking the circularity.

\smallskip
\noindent\textbf{Future direction: LLM-assisted annotation suggestion and auto-repair.}
A natural way to lower both barriers is to combine IntentChecker with a large language model that observes the same source the developer is editing.
For intent annotations, the LLM could generate candidate @During/@Post annotations by reading the surrounding function, the contract's NatSpec comments, and the developer's recent commit messages, presenting the candidates as suggestions that the developer accepts or edits.
For debug annotations, the LLM could propose representative ranges based on the function's call sites elsewhere in the codebase or based on standard token decimals (e.g., 18 for most ERC20 tokens, 6 for USDC).
Beyond annotation suggestion, an LLM that has access to IntentChecker's Violated judgment and violated region could also propose code edits that move the violated region out of the reachable interval, effectively turning the awareness signal into a candidate fix.
We deliberately leave both directions as future work; what matters for this paper is that they are natural extensions enabled by, rather than blocked by, the current design.

\smallskip
\noindent\textbf{Developer study.}
This paper evaluates IntentChecker through structural analysis---what the tool can and cannot detect (RQ1), what the intent model can and cannot express (RQ2), and how the tool compares with existing detectors (RQ3)---rather than through a developer study measuring perceived usability or annotation intuitiveness.
A controlled study with Solidity-experienced participants would answer a complementary set of questions: whether developers find the annotation syntax natural, how long annotation authoring takes in practice, and whether the tool's feedback leads to faster bug discovery compared to existing debugging workflows.
Designing such a study requires careful attention to participant recruitment (sufficient Solidity expertise) and ecological validity (realistic development tasks rather than isolated code snippets), and we defer it to future work as an independent evaluation axis that builds on the structural foundation established here.

\subsection{Remaining Technical Limitations}\label{sec:disc:tech}

In addition to the annotation authoring burden discussed above, three classes of technical limitations remain in the current implementation.

\smallskip
\noindent\textbf{Multi-transaction invariants.}
IntentChecker's analysis scope is multi-contract but single-transaction.
A whole class of important numeric properties cannot be expressed within a single transaction because they relate quantities established by different transactions in the contract's lifetime.
The canonical example is the ERC20 conservation invariant totalSupply == sum(balances): any single function execution may temporarily increment one side and not the other, but across the entire transaction history the two must agree.
Another example is the price oracle manipulation pattern, where the buggy behavior only manifests when one transaction sets up state that a later transaction reads.
Validating such properties would require tracking state across transactions and aggregating across all balance entries, both of which lie outside the current single-transaction interval-domain analysis.
This is a known boundary of the development-time intent validation approach we have taken, and is more naturally addressed by post-hoc tools or formal verification frameworks that analyze entire contract lifecycles.

\smallskip
\noindent\textbf{Non-arithmetic Yul and other unsupported Solidity constructs.}
The arithmetic-focused Yul subset described in §\ref{sec:prelim:extension}(c) covers the cases where inline assembly is used as a faster substitute for high-precision integer arithmetic, which is by far the most common use case in DeFi libraries.
However, hand-written Yul that uses memory and storage opcodes (mload/mstore, sload/sstore), control flow (if, for, switch), or function definitions falls outside the supported subset and is conservatively treated as $\top$.
The same conservative treatment applies to a small number of Solidity constructs that the abstract interpreter does not model: abi.decode from raw bytes, keccak256, address(this).balance (whose value depends on the EVM-level execution context rather than on tracked state), and other built-ins whose results depend on opaque cryptographic or VM-level operations.
Bugs whose buggy/correct distinction depends on the precise value produced by such constructs are correspondingly out of IntentChecker's reach (these account for the L3 cases in Figure~\ref{fig:rq1-limitations}).
Extending the Yul visitor to handle memory and storage opcodes and the abstract interpreter to handle a wider set of opaque builtins is a meaningful but bounded engineering task.

\smallskip
\noindent\textbf{Soundness bounded by debug annotation scope.}
IntentChecker's analysis is sound within the ranges declared by debug annotations: when the validation engine reports Satisfied, the declared intent provably holds for every combination of values in those ranges---this follows from the set-theoretic definition of the comparison primitive (§\ref{sec:intent-semantics}) applied pointwise to each annotation.
The guarantee does not extend outside those ranges.
A bug that only manifests at the boundary of a parameter's representable values, but not in the range the developer chose for @LocalVar, will not appear in any judgment.
This is an inherent characteristic of debugging tools---traditional debuggers also only execute the inputs the developer provides---and we encourage developers to use representative ranges and to iteratively refine them when Warning results draw attention to specific sub-intervals.

\subsection{Practical Annotation Guidelines}\label{sec:disc:guidelines}

The annotation grammar (§\ref{sec:design}) defines what can be written; the question that remains is when a developer should write which annotation.
The 20 mitigated cases and 14 L5 expressible cases reveal a recurring pattern: the most impactful annotations are not the ones that verify what the developer did write, but the ones that catch what the developer forgot to write.
We distill three non-obvious guidelines from this observation.

\smallskip
\noindent\textbf{Guideline 1: Enumerate all state variables a function should affect, not just those it does affect.}
Developers naturally reason about the code they wrote---the variables they explicitly assigned, the calls they explicitly made.
Bugs of omission (L5a, 7 cases) arise precisely from this cognitive bias: a state variable that should have been updated was simply never mentioned.
After writing a function, the developer should list every state variable whose value the function is expected to change and add at least a directional annotation (@Post var(Entry\,$<$\,Exit) or @Post changed(var, true)) for each.
In our data, this single practice---checking what should change rather than what does change---would have caught 7 of the 14 L5 expressible cases and 5 of the 8 mitigated @Post cases, all without requiring any knowledge of the specific bug.

\smallskip
\noindent\textbf{Guideline 2: When two state variables should move together, annotate their relationship, not each one in isolation.}
Several benchmark cases involve accounting pairs (e.g., a per-user balance and a global total, a deposit ledger and a reward tracker) that should stay in sync.
Annotating each variable independently catches only missing-update bugs; annotating their relationship (e.g., @Post globalTotal(Entry\,-\,Exit) $\geq$ userBalance(Entry\,-\,Exit)) catches bugs where both variables are updated but by inconsistent amounts.
This pattern recurs in 4 of the 20 mitigated cases (all Erroneous Accounting) and is the dominant form in L4d (invariant-masked), where independent annotations on each variable succeed but the cross-variable invariant fails.

\smallskip
\noindent\textbf{Guideline 3: For multi-step arithmetic, annotate the result against the mathematically ideal expression, not the code as written.}
Solidity's integer division truncates at every intermediate step, and the resulting precision loss depends on the order of operations.
Code that reads amount.div(12).mul(5) may look correct, but the developer's actual intent is the mathematical value $\lfloor\text{amount} \times 5 / 12\rfloor$.
Writing the intent annotation in the ideal form---@During result $\geq$ amount * 5 / 12---causes IntentChecker to compare the truncated code path against the non-truncated reference, surfacing the discrepancy automatically.
This pattern accounts for 7 of the 20 mitigated cases (Div In Path, Operator Order Issue, Exchange Problem, and Precision Loss Trend), all of which were caught by annotations that express the mathematically correct expression rather than restating what the code already does.

\section{Related Works}\label{sec:related}
\subsection{Numeric Logic Error Detection in Solidity}
Logic error detection focuses on identifying discrepancies between developer intentions and actual program behavior in smart contracts.

\citep{highguard} addresses business logic violations in smart contracts by using Dynamic Condition Response (DCR) graphs
to model expected behaviors including temporal constraints, role-based access control, and inter-function dependencies.
The system monitors runtime transactions against user-provided DCR models and their function mappings,
generating violation reports when execution activities deviate from the specified business logic patterns.

\citep{smartpulse} verifies temporal properties of smart contracts by accepting user-specified Linear Temporal Logic (LTL) formulas
and attack models to check both safety and liveness properties.
The system not only detects violations but also generates concrete attack scenarios that demonstrate how the temporal properties can be violated,
providing developers with actionable insights into potential logic errors.

\citep{consolidating} enables developers to specify function and address behaviors using the ConSol syntax alongside their Solidity code,
then verifies whether the implementation adheres to these behavioral specifications and generates behavior-compliant Solidity programs without annotations.

\citep{smartinv} leverages multimodal learning to automatically infer invariants and detect machine unauditable bugs (MUBs)~\citep{exploitablebugs}
by analyzing both source code and natural language artifacts including comments, function signatures, and documented transaction scenarios.
The system employs a fine-tuned model with Tier of Thought (ToT) prompting that progressively predicts transaction contexts,
generates and ranks invariants, and ultimately identifies vulnerabilities through three-tier reasoning.

\citep{accountingerrors} proposes a novel type system for Solidity variables defined as tuples of $\langle$Financial Type, Scaling Factor, Token Unit, Address$\rangle$
to detect accounting errors through type checking.
The system reinterprets Solidity code based on detection patterns such as incorrectly adding balances with fees,
enabling systematic identification of numeric logic errors that arise from mismatched financial operations.

Furthermore, there are also approaches that address multiple transaction logic errors such as price manipulation~\citep{defiranger,defiguard,pomabuster,secplf,defitainter}
and flash loan attacks~\citep{defitail,midnight}.

The numeric logic error detectors above and IntentChecker occupy distinct points along the development--auditing timeline rather than competing on the same axis.
NumScout, GPTScan, SmartInv, and the type-system approach of \citet{accountingerrors} are all designed to run over completed code and recover developer intent indirectly through pattern catalogs, language-model inference, or financial type assignments.
IntentChecker is positioned earlier in time and lets the developer declare intent explicitly through annotations evaluated inside the editor.
Because the two approaches see different inputs (code only vs.\ code with declared intent) and run in different workflow stages (post-hoc audit vs.\ development-time validation), their coverage is necessarily complementary, and §\ref{sec:rq3} characterizes that complementarity rather than arguing numerical superiority.

\subsection{Specification-based Smart Contract Verification}
Specification-based verification approaches focus on verifying general properties and correctness requirements of smart contracts rather than specific developer intentions.

\citep{scribble} provides a framework that transforms developer annotations into Solidity assertion code,
enabling testing with existing tools like MythX~\citep{mythx} and Mythril~\citep{mythril}.
The framework supports multiple annotation types including if\_succeed (conditions for successful function returns),
invariant (contract-level properties that must always hold), if\_updated (conditions checked when state variables change),
and assert (validations at specific execution points), making formal specifications accessible through familiar testing pipelines.

\citep{smartace} builds upon Scribble~\citep{scribble} to verify whether Solidity code satisfies specifications written in the Scribble language,
leveraging C-based verification tools like SeaHorn~\citep{seahorn}, libFuzzer~\citep{libfuzzer}, and KLEE~\citep{klee} by transforming smart contracts into LLVM-IR based harnesses.
The framework abstracts the $2^{160}$ possible user addresses into a small set of representative users (implicit, transient, and explicit),
enabling verification of properties that must hold for arbitrary numbers of users, such as accounting properties like
"the contract balance equals the sum of all user deposits."

\citep{idealist} verifies Solidity code using two types of annotations: specification annotations (loop invariants, call invariants, and call modifies) that developers declare as assumptions,
and verification annotations (ensures, reverts\_if, and contract invariants) that the tool checks against the implementation.
The system assumes the truth of specification annotations and uses them to verify whether the code satisfies function post-conditions,
correctly handles revert conditions, and maintains contract invariants throughout execution.

\citep{solvent} focuses on verifying liquidity properties of smart contracts through formal methods
by converting both the contract code and user-defined liquidity properties into SMT constraints.
The system then checks whether the contract satisfies these properties using SMT solvers,
generating counterexamples when violations are detected to help developers understand potential liquidity issues.


While specification-based approaches formalise contract-level properties in dedicated specification languages---Scribble's annotation language, Idealist's loop and call invariants, Solvent's SMT-encoded liquidity properties---and discharge verification obligations over completed code, IntentChecker's intent model occupies a different point in the design space with three distinguishing characteristics.
First, its annotation language targets value-level numeric intent on individual variables at specific program points, not contract-level invariants or temporal properties.
Second, its denotational semantics (§\ref{sec:intent-semantics}) maps each annotation form to a tri-valued validation outcome over the abstract states produced by interval-domain analysis, giving the annotation language a formal foundation comparable in rigor to the specification languages above while remaining cheap enough to evaluate at development time.
Third, because IntentChecker performs interval-domain analysis that tracks actual value ranges through library function calls (e.g., SafeMath operations) rather than treating libraries as black boxes, it can validate intent over step-by-step numeric computations rather than only over post-condition-style abstract properties.
The two approaches are therefore complementary in workflow stage (development-time vs.\ audit/verification-time), in granularity (value-level vs.\ contract-level), and in the specification style they support.

\subsection{Tools for Strengthening Security During Development Phase}
Development-phase security tools integrate vulnerability detection directly into the coding workflow, enabling developers to identify and address security issues as they write code rather than in post-development analysis.
\citep{making} presents an integrated development environment for smart contract development that combines natural language model-based code autocompletion with security validation.
The platform integrates existing static analysis tools including Oyente~\citep{oyente}, Mythril~\citep{mythril}, and Securify~\citep{securify} to provide comprehensive security reports during development workflow.

\citep{solhint} provides a Solidity linter that detects predefined vulnerability patterns including low-level calls, reentrancy, and deprecated function usage during development.
The tool enforces security best practices by checking for patterns such as proper visibility declarations, send result validation, and secure compiler version usage.

While these development-phase tools rely on predefined detection patterns for known vulnerabilities,
IntentChecker enables developers to specify their intended numeric behaviors through annotations and validates whether the actual program execution matches these specified intentions.
IntentChecker also extends SolQDebug's analysis scope beyond a single contract to a multi-contract single-transaction scope that includes inheritance, libraries, interfaces, and an arithmetic Yul subset, which is what makes intent validation feasible on realistic DeFi code.

\section{Conclusion}\label{sec:conclusion}

We presented IntentChecker, a development-time intent validation tool that extends SolQDebug with (i) an intent model---two annotation types (@During/@Post) equipped with a formal denotational semantics that maps each annotation to a tri-valued outcome against abstract states, (ii) a validation engine whose interval comparison primitive is proved sound against that semantics, producing graded Satisfied/Warning/Violated signals with separately computed risk scores and violated regions, and (iii) a scope extension to multi-contract single-transaction analysis covering inheritance, libraries, interfaces, and an arithmetic Yul subset.

On 75 effective benchmark instances (89 collected, 14 excluded) from Web3Bugs and NumScout, IntentChecker surfaced 20 bugs via Violated judgments at interactive speeds (mean 3.69\,s per case).
RQ2's expressibility analysis showed that 14 of the 34 intent-model limitation cases are expressible under bug awareness, while 20 represent a hard grammar ceiling---identifying concrete extension points for the intent model.
IntentChecker's positioning relative to existing detection tools is complementary: it fills the development-time gap that auditing-time tools structurally cannot occupy.

The primary practical limitation is the annotation authoring burden, which future work can address through LLM-assisted annotation suggestion (lowering the L5 bug-awareness ceiling) and auto-repair from violated regions.
On the analyzer side, extending Yul support and multi-transaction reasoning would close the remaining technical limitations.
More broadly, this work adds ``development-time intent validator'' as a third category alongside pattern-based detectors and formal verifiers in the design space for numeric logic error tools.

\section*{Acknowledgements}
This work was supported by the Institute of Information \& communications Technology Planning \& Evaluation (IITP) grant funded by the Korean government (MSIT) (RS-2021-II210177, High Assurance of Smart Contract for Secure Software Development Life Cycle), the Technology development Program (RS-2025-25455127) funded by the Ministry of SMEs and Startups (MSS, Korea), and the ``Regional Innovation System \& Education (RISE)'' through the Seoul RISE Center, funded by the Ministry of Education (MOE) and the Seoul Metropolitan Government (2025-RISE-01-003-01).

\section*{Author Contributions}
Inseong Jeon participated in conceptualization, methodology design, system implementation, data collection, experiments, and manuscript writing. Sundeuk Kim and Hyunwoo Kim assisted in experiments, data collection and analysis, and contributed to manuscript writing. Hoh Peter In provided resources, assisted in editing the manuscript, and supervised the entire project. All authors reviewed and approved the final version of the manuscript.

\section*{Data Availability}
The benchmark of 89 real-world numeric logic error instances curated from Web3Bugs~\citep{exploitablebugs}, NumScout~\citep{numscout}, and the Greedy Contract dataset of \citet{flyinointment}, along with the IntentChecker implementation~\citep{intentchecker}, evaluation scripts, and experimental results, are available at \url{https://github.com/iwwyou/IntentChecker.git}.

\section*{Declarations}

\noindent\textbf{Competing interests} The authors declare no competing interests.

\noindent\textbf{Ethical approval} Not applicable since there are no human and/or animal studies included in this paper.


\begin{thebibliography}{99}

\bibitem[Arora et~al.(2024)]{secplf}
Arora, S., et al.: SecPLF: Secure protocols for loanable funds against oracle manipulation attacks. In: Proceedings of the 19th ACM Asia Conference on Computer and Communications Security (ASIACCS), pp. 1394--1405 (2024). https://doi.org/10.1145/3634737.3637681

\bibitem[Bartoletti et~al.(2024)]{solvent}
Bartoletti, M., et al.: Solvent: liquidity verification of smart contracts. In: International Conference on Integrated Formal Methods (iFM), pp. 256--266 (2024). \url{https://doi.org/10.1007/978-3-031-76554-4_14}

\bibitem[Cadar et~al.(2008)]{klee}
Cadar, C., et al.: KLEE: unassisted and automatic generation of high-coverage tests for complex systems programs. In: OSDI, pp. 209--224 (2008). https://doi.org/10.5555/1855741.1855756

\bibitem[Chaliasos et~al.(2024)]{defisecurity}
Chaliasos, S., et al.: Smart contract and defi security tools: Do they meet the needs of practitioners?. In: Proceedings of the 46th IEEE/ACM International Conference on Software Engineering (ICSE), pp. 1--13 (2024). https://doi.org/10.1145/3597503.3623302

\bibitem[Chen et~al.(2025)]{numscout}
Chen, J., et al.: NumScout: Unveiling Numerical Defects in Smart Contracts using LLM-Pruning Symbolic Execution. IEEE Transactions on Software Engineering (2025). https://doi.org/10.1109/TSE.2025.3555622

\bibitem[NumScout(2026)]{numscout_dataset}
NumScout: 95 Samples Experiment Dataset. \url{https://github.com/NumScout/NumScout/tree/main/Experiment/95_Samples_Run} (2026). Accessed April 2026

\bibitem[Consensys Diligence(2025a)]{mythril}
Consensys Diligence: Mythril. \url{https://github.com/ConsenSysDiligence/mythril} (2025). Accessed December 2025

\bibitem[Consensys Diligence(2025b)]{mythx}
Consensys Diligence: MythX. \url{https://mythx.io/} (2025). Accessed December 2025

\bibitem[Consensys Diligence(2025c)]{scribble}
Consensys Diligence: Scribble. \url{https://diligence.security/scribble/} (2025). Accessed January 2026

\bibitem[Eshghie et~al.(2024)]{highguard}
Eshghie, M., et al.: Highguard: Cross-chain business logic monitoring of smart contracts. In: Proceedings of the 39th IEEE/ACM International Conference on Automated Software Engineering (ASE), pp. 2378--2381 (2024). https://doi.org/10.1145/3691620.3695356

\bibitem[Feist et~al.(2019)]{slither}
Feist, J., Grieco, G., Groce, A.: Slither: A static analysis framework for smart contracts. In: 2019 IEEE/ACM 2nd International Workshop on Emerging Trends in Software Engineering for Blockchain (WETSEB), pp. 8--15 (2019). https://doi.org/10.1109/WETSEB.2019.00008

\bibitem[Gurfinkel et~al.(2015)]{seahorn}
Gurfinkel, A., et al.: The SeaHorn verification framework. In: International Conference on Computer Aided Verification, pp. 343--361 (2015). \url{https://doi.org/10.1007/978-3-319-21690-4_20}

\bibitem[Jeon et~al.(2025)]{solqdebug}
Jeon, I., et al.: SolQDebug: Debug Solidity Quickly for Interactive Immediacy in Smart Contract Development (2025). https://doi.org/10.21203/rs.3.rs-8077153/v1

\bibitem[IntentChecker(2026a)]{intentchecker}
IntentChecker. \url{https://github.com/iwwyou/IntentChecker.git} (2026). Accessed April 2026

\bibitem[IntentChecker(2026b)]{dappscananalysis}
IntentChecker: DAppSCAN contract analysis dataset. \url{https://github.com/iwwyou/IntentChecker/tree/main/evaluation/DappSCAN_contract_analysis} (2026). Accessed April 2026

\bibitem[Kong et~al.(2023)]{defitainter}
Kong, Q., et al.: Defitainter: Detecting price manipulation vulnerabilities in defi protocols. In: Proceedings of the 32nd ACM SIGSOFT International Symposium on Software Testing and Analysis (ISSTA), pp. 1144--1156 (2023). https://doi.org/10.1145/3597926.3598124

\bibitem[Li et~al.(2025)]{defitail}
Li, X., et al.: Penetrating the hostile: Detecting defi protocol exploits through cross-contract analysis. IEEE Transactions on Information Forensics and Security 20, 11759--11774 (2025). https://doi.org/10.1109/TIFS.2025.3626979

\bibitem[LLVM(2025)]{libfuzzer}
LLVM: libFuzzer. \url{https://llvm.org/docs/LibFuzzer.html} (2025). Accessed December 2025

\bibitem[Luu et~al.(2016)]{oyente}
Luu, L., et al.: Making smart contracts smarter. In: Proceedings of the 2016 ACM SIGSAC Conference on Computer and Communications Security (CCS), pp. 254--269 (2016). https://doi.org/10.1145/2976749.2978309

\bibitem[Nguyen et~al.(2023)]{idealist}
Nguyen, T.D., et al.: An idealist's approach for smart contract correctness. In: International Conference on Formal Engineering Methods, pp. 11--28 (2023). \url{https://doi.org/10.1007/978-981-99-7584-6_2}

\bibitem[Pearson(2025)]{pearson_r}
Pearson correlation coefficient. \url{https://en.wikipedia.org/wiki/Pearson_correlation_coefficient} (2025). Accessed April 2026

\bibitem[Protofire(2025)]{solhint}
Protofire: Solhint. \url{https://github.com/protofire/solhint} (2025). Accessed Janauary 2026

\bibitem[Ramezany et~al.(2023)]{midnight}
Ramezany, S., Setthawong, R., Setthawong, P.: Midnight: An Efficient Event-driven EVM Transaction Security Monitoring Approach For Flash Loan Detection. In: Proceedings of the 2023 20th International Joint Conference on Computer Science and Software Engineering (JCSSE), pp. 237--241 (2023). https://doi.org/10.1109/JCSSE58229.2023.10201970

\bibitem[Ren et~al.(2021)]{making}
Ren, M., et al.: Making smart contract development more secure and easier. In: Proceedings of the 29th ACM Joint Meeting on European Software Engineering Conference and Symposium on the Foundations of Software Engineering (ESEC/FSE), pp. 1360--1370 (2021). https://doi.org/10.1145/3468264.3473929

\bibitem[Soud et~al.(2024)]{flyinointment}
Soud, M., Liebel, G., Hamdaqa, M.: A fly in the ointment: an empirical study on the characteristics of Ethereum smart contract code weaknesses. Empirical Software Engineering 29(1), 13 (2024). https://doi.org/10.1007/s10664-023-10398-5

\bibitem[Stephens et~al.(2021)]{smartpulse}
Stephens, J., et al.: SmartPulse: automated checking of temporal properties in smart contracts. In: Proceedings of the 2021 IEEE Symposium on Security and Privacy (SP), pp. 555--571 (2021). https://doi.org/10.1109/SP40001.2021.00085

\bibitem[Sun et~al.(2024)]{gptscan}
Sun, Y., et al.: Gptscan: Detecting logic vulnerabilities in smart contracts by combining gpt with program analysis. In: Proceedings of the IEEE/ACM 46th International Conference on Software Engineering (ICSE), pp. 1--13 (2024). https://doi.org/10.1145/3597503.3639117

\bibitem[Tsankov et~al.(2018)]{securify}
Tsankov, P., et al.: Securify: Practical security analysis of smart contracts. In: Proceedings of the 2018 ACM SIGSAC Conference on Computer and Communications Security (CCS), pp. 67--82 (2018). https://doi.org/10.1145/3243734.3243780

\bibitem[Wang et~al.(2024a)]{defiguard}
Wang, D., et al.: Defiguard: A price manipulation detection service in defi using graph neural networks. IEEE Transactions on Services Computing (2024). https://doi.org/10.1109/TSC.2024.3489439

\bibitem[Wang et~al.(2024b)]{smartinv}
Wang, S.J., Pei, K., Yang, J.: Smartinv: Multimodal learning for smart contract invariant inference. In: Proceedings of the 2024 IEEE Symposium on Security and Privacy (SP), pp. 2217--2235 (2024). https://doi.org/10.1109/SP54263.2024.00126

\bibitem[Wei et~al.(2024)]{consolidating}
Wei, G., et al.: Consolidating Smart Contracts with Behavioral Contracts. Proceedings of the ACM on Programming Languages 8(PLDI), 965--989 (2024). https://doi.org/10.1145/3656416

\bibitem[Wesley et~al.(2022)]{smartace}
Wesley, S., et al.: Verifying Solidity smart contracts via communication abstraction in SmartACE. In: International Conference on Verification, Model Checking, and Abstract Interpretation, pp. 425--449 (2022). \url{https://doi.org/10.1007/978-3-030-94583-1_21}

\bibitem[Wu et~al.(2023)]{defiranger}
Wu, S., et al.: Defiranger: Detecting defi price manipulation attacks. IEEE Transactions on Dependable and Secure Computing 21(4), 4147--4161 (2023). https://doi.org/10.1109/TDSC.2023.3346888

\bibitem[Xi et~al.(2024)]{pomabuster}
Xi, R., Wang, Z., Pattabiraman, K.: POMABuster: Detecting Price Oracle Manipulation Attacks in Decentralized Finance. In: Proceedings of the 2024 IEEE Symposium on Security and Privacy (SP), pp. 3923--3942 (2024). https://doi.org/10.1109/SP54263.2024.00257

\bibitem[Zhang(2024)]{accountingerrors}
Zhang, B.: Towards finding accounting errors in smart contracts. In: Proceedings of the IEEE/ACM 46th International Conference on Software Engineering (ICSE), pp. 1--13 (2024). https://doi.org/10.1145/3597503.3639128

\bibitem[Zhang et~al.(2024)]{nyx}
Zhang, W., et al.: Nyx: Detecting exploitable front-running vulnerabilities in smart contracts. In: Proceedings of the 2024 IEEE Symposium on Security and Privacy (SP), pp. 2198--2216 (2024). https://doi.org/10.1109/SP54263.2024.00146

\bibitem[Yang et~al.(2025)]{hyperion}
Yang, S., et al.: Hyperion: Unveiling DApp Inconsistencies Using LLM and Dataflow-Guided Symbolic Execution. In: Proceedings of the 2025 IEEE/ACM 47th International Conference on Software Engineering (ICSE), pp. 2125--2137 (2025). https://doi.org/10.1109/ICSE55347.2025.00015

\bibitem[Zhang et~al.(2023)]{exploitablebugs}
Zhang, Z., et al.: Demystifying exploitable bugs in smart contracts. In: Proceedings of the 2023 IEEE/ACM 45th International Conference on Software Engineering (ICSE), pp. 615--627 (2023). https://doi.org/10.1109/ICSE48619.2023.00061

\bibitem[Zheng et~al.(2024)]{dappscan}
Zheng, Z., et al.: Dappscan: building large-scale datasets for smart contract weaknesses in dApp projects. IEEE Transactions on Software Engineering 50(6), 1360--1373 (2024). https://doi.org/10.1109/TSE.2024.3383422

\bibitem[Zhou et~al.(2023)]{defiattacks}
Zhou, L., et al.: Sok: Decentralized finance (defi) attacks. In: Proceedings of the 2023 IEEE Symposium on Security and Privacy (SP), pp. 2444--2461 (2023). https://doi.org/10.1109/SP46215.2023.10179435

\end{thebibliography}


\end{document}
